% !TEX program = lualatex
% Options for packages loaded elsewhere
\PassOptionsToPackage{unicode}{hyperref}
\PassOptionsToPackage{hyphens}{url}
\documentclass[
]{article}
\usepackage{xcolor}
\usepackage{graphicx}
\usepackage{amsmath,amssymb}
\setcounter{secnumdepth}{-\maxdimen} % remove section numbering
\usepackage{iftex}
\ifPDFTeX
  \usepackage[T1]{fontenc}
  \usepackage[utf8]{inputenc}
  \usepackage{textcomp} % provide euro and other symbols
\else % if luatex or xetex
  \usepackage{unicode-math} % this also loads fontspec
  \defaultfontfeatures{Scale=MatchLowercase}
  \defaultfontfeatures[\rmfamily]{Ligatures=TeX,Scale=1}
\fi
\usepackage{lmodern}
\ifPDFTeX\else
  % xetex/luatex font selection
\fi
% Use upquote if available, for straight quotes in verbatim environments
\IfFileExists{upquote.sty}{\usepackage{upquote}}{}
\IfFileExists{microtype.sty}{% use microtype if available
  \usepackage[]{microtype}
  \UseMicrotypeSet[protrusion]{basicmath} % disable protrusion for tt fonts
}{}
\makeatletter
\@ifundefined{KOMAClassName}{% if non-KOMA class
  \IfFileExists{parskip.sty}{%
    \usepackage{parskip}
  }{% else
    \setlength{\parindent}{0pt}
    \setlength{\parskip}{6pt plus 2pt minus 1pt}}
}{% if KOMA class
  \KOMAoptions{parskip=half}}
\makeatother
\setlength{\emergencystretch}{3em} % prevent overfull lines
\providecommand{\tightlist}{%
  \setlength{\itemsep}{0pt}\setlength{\parskip}{0pt}}
\usepackage{bookmark}
\IfFileExists{xurl.sty}{\usepackage{xurl}}{} % add URL line breaks if available
\urlstyle{same}
\hypersetup{
  pdftitle={Pivato Paper},
  hidelinks,
  pdfcreator={LaTeX via pandoc}}

\title{Pivato Paper}
\author{}
\date{}

\begin{document}
\maketitle

\protect\phantomsection\label{s000005}
1Introduction

Suppose a group of voters must collectively choose some policy from a
set of alternatives, using a voting rule \(F\). Suppose we split the
voters into two subgroups, and each subgroup, using rule \(F\), selects
the alternative \(x\). Then it seems desirable that the combined group,
using \(F\), should also select alternative \(x\). We say the rule \(F\)
satisfies reinforcement\textsuperscript{1}1Smith (1973) calls this
condition `separability', while Young, 1974b, Young, 1975 and~Saari
(1990) call it `consistency'. Dhillon (1998) and Dhillon and Mertens
(1999) refer to an analogous property as `extended Pareto', while Gilboa
and Schmeidler (2003) call it `combination'. if it has this property.
Smith (1973) and Young, 1974b, Young, 1975 showed that `scoring rules'
are the only preference aggregation rules\textsuperscript{2}2Smith
(1973) and Young (1974b) consider rules which produce social preference
relations as output, whereas Young (1975) considers rules which produce
subsets of social alternatives as output; these two frameworks yield two
slightly different definitions of `reinforcement' and `scoring rule'.
which satisfy reinforcement and are anonymous and neutral
(i.e.~invariant under relabeling of the voters and/or alternatives).
These results have led to characterizations of the Borda
rule,\textsuperscript{3}3See~Young, 1974a, Young, 1975 and Nitzan and
Rubinstein (1981). Kemeny rule,\textsuperscript{4}4See Young and
Levenglick (1978). and plurality rule\textsuperscript{5}5See Richelson
(1978), Morkelyunas (1982), Ching (1996), and~Yeh (2008). each as the
only preference aggregation rule which satisfies reinforcement along
with certain other axioms.

Myerson (1995) generalized the Smith--Young results from preference
aggregators to abstract voting rules. In Myerson's framework, there is a
finite set of social alternatives \(\mathcal{X}\) and a finite set of
`signals' \(\mathcal{V}\). A `profile' assigns a signal to each voter,
and an abstract voting rule selects some nonempty subset of
\(\mathcal{X}\) for each profile. A scoring rule is a voting rule where
each element of \(\mathcal{V}\) assigns a real-valued `score' to each
element of \(\mathcal{X}\); the rule then selects the alternative(s)
with the highest total score.\textsuperscript{6}6Zwicker (2008) has
shown that such scoring rules can also be interpreted as `mean proximity
rules': each element of \(\mathcal{V}\) and \(\mathcal{X}\) is
represented as a vector in \(\mathbb{R}^{N}\), and the social choice is
the element of \(\mathcal{X}\) which is closest to the vector average of
the signals of the voters. Myerson showed that if an abstract voting
rule satisfied reinforcement, universal domain (i.e.~it is defined for
all profiles), a form of neutrality (i.e.~all social alternatives are
treated equally), and an Archimedean/continuity condition he called
overwhelming majority, then it was a scoring rule. For example, approval
voting (Brams and Fishburn, 1983) is a scoring rule---indeed, it is the
only abstract voting rule which satisfies reinforcement along with
certain other axioms.\textsuperscript{7}7See Fishburn (1978),
Morkelyunas (1981), and Alós-Ferrer (2006).

In Section~2, I extend Myerson's representation theorem, by considering
infinite signal sets, and removing the hypotheses of universal domain
and overwhelming majority. I do this by considering scoring rules where
the scores can range over an abstract linearly ordered abelian group,
instead of ranging over the real numbers.\textsuperscript{8}8Smith
(1973) also considered linearly ordered abelian groups. I also
characterize balance rules (another class introduced by Myerson, 1995).
Next, Section~3 shows that the `neutrality'
(i.e.~permutation-equivariance) properties of a voting rule can be
reflected by corresponding neutrality properties of its balance
representation or scoring representation. Section~4 concludes with some
open problems. Appendix~A gives the technical arguments underlying some
of the examples which appear in the text. Appendix~B provides some
background on linearly ordered abelian groups. Appendix~C contains the
proofs of all results in the paper.

\protect\phantomsection\label{s000010}
2Model and main results

Let \(\mathcal{X}\) be a set of social alternatives, and let
\(\mathcal{V}\) be the set of possible signals which could be sent by
each voter. (The sets \(\mathcal{X}\) and \(\mathcal{V}\) could be
finite or infinite.)~Let
\(\mathbb{N} ≔ \left\{ 0,1,2,3,\ldots \right\}\) and
\(\mathbb{Z} ≔ \left\{ \pm n;n \in \mathbb{N} \right\}\). For any
\(\mathbf{n} \in \mathbb{Z}^{\mathcal{V}}\), let
\(\left\| \mathbf{n} \right\|_{} ≔ \sum_{v \in \mathcal{V}}\left| n_{v} \right|\).
Define
\(\mathbb{N}^{\langle\mathcal{V}\rangle} ≔ \left\{ \mathbf{n} \in \mathbb{N}^{\mathcal{V}};\left\| \mathbf{n} \right\|_{} < \infty \right\}\).
If \(\mathbf{n} \in \mathbb{N}^{\langle\mathcal{V}\rangle}\), then
\(\mathbf{n}\) represents an anonymous profile of voters: for each
\(v \in \mathcal{V}\), we interpret \(n_{v}\) as the number of voters
sending the signal \(v\), while \(\left\| \mathbf{n} \right\|_{}\) is
the size of the whole population. Note that we do not fix
\(\left\| \mathbf{n} \right\|_{}\) in advance. Let
\(\mathbf{0} = (0,0,\ldots,0)\) be the all-zeros vector in
\(\mathbb{N}^{\langle\mathcal{V}\rangle}\). A domain is any collection
of profiles
\(\mathcal{D} \subseteq \mathbb{N}^{\langle\mathcal{V}\rangle}\) such
that \(\mathbf{0} \in \mathcal{D}\). (The set
\(\mathbb{N}^{\langle\mathcal{V}\rangle}\) itself is the universal
domain.)~A variable population, anonymous voting rule is a
correspondence
\(\left. F:\mathcal{D}\rightrightarrows\mathcal{X} \right.\) such that
\(F\left( \mathbf{0} \right) = \mathcal{X}\) (this property is often
called the axiom of general abstention). Thus, for all
\(\mathbf{d} \in \mathcal{D}\), the outcome
\(F\left( \mathbf{d} \right) \subseteq \mathcal{X}\) is a nonempty set
(typically a singleton).

A linearly ordered abelian group is a triple
\(\left( \mathcal{R}, + , > \right)\), where \(\mathcal{R}\) is a set,
`\(+\)' is an abelian group operation, and `\(>\)' is a complete,
antisymmetric, transitive binary relation such that, for all
\(r,s \in \mathcal{R}\), if \(r > 0\), then \(r + s > s\). For example:
the set \(\mathbb{R}\) of real numbers is a linearly ordered abelian
group, with the standard ordering and addition operator. So is any
subgroup of \(\mathbb{R}\). For any nonzero \(N \in \mathbb{N}\), let
\(\mathbb{R}_{lex}^{N}\) denote the group \(\mathbb{R}^{N}\) with vector
addition and the lexicographic order \(\gg\)\textsuperscript{9}9To be
precise: for any \(\mathbf{r} = \left( r_{1},\ldots,r_{N} \right)\) and
\(\mathbf{r}' ≔ \left( r_{1}',\ldots,r_{N}' \right)\) in
\(\mathbb{R}^{N}\), we have \(\mathbf{r} \gg \mathbf{r}'\) if there is
some \(M \in \lbrack 1\ldots N\rbrack\) such that \(r_{n} = r_{n}'\) for
all \(n < M\), while \(r_{M} > r_{M}'\).; then \(\mathbb{R}_{lex}^{N}\)
is a linearly ordered abelian group. Such groups are useful for
representing preferences where some decision variables have
lexicographical priority over others.

For any
\(\mathbf{r} = \left( r_{v} \right)_{v \in \mathcal{V}} \in \mathcal{R}^{\mathcal{V}}\),
we define a function
\(\left. \mathbf{r}:\mathbb{N}^{\langle\mathcal{V}\rangle}\longrightarrow\mathcal{R} \right.\)
by setting
\(\mathbf{r}\left( \mathbf{n} \right) ≔ \sum_{v \in \mathcal{V}}n_{v}r_{v}\)
for all
\(\mathbf{n} \in \mathbb{N}^{\langle\mathcal{V}\rangle}\).\textsuperscript{10}10Thus,
if \(\mathcal{R} = \mathbb{R}\), then
\(\mathbf{r}\left( \mathbf{n} \right) = \mathbf{r}\bullet\mathbf{n}\),
where `\(\bullet\)' is the inner product operation on
\(\mathbb{R}^{\mathcal{V}}\). An \(\mathcal{R}\)-valued score system on
\(\left( \mathcal{X},\mathcal{V} \right)\) is an \(\mathcal{X}\)-indexed
collection
\(\mathsf{S} ≔ \left\{ \mathbf{s}^{x} \right\}_{x \in \mathcal{X}} \subset \mathcal{R}^{\mathcal{V}}\).
For any domain
\(\mathcal{D} \subseteq \mathbb{N}^{\langle\mathcal{V}\rangle}\), the
scoring rule determined by \(\mathsf{S}\) is the voting rule
\(\left. F_{\mathsf{S}}:\mathcal{D}\rightrightarrows\mathcal{X} \right.\)
defined as follows:
\[F_{\mathsf{S}}\left( \mathbf{d} \right) ≔ \underset{x \in \mathcal{X}}{arg~max}\mathbf{s}^{x}\left( \mathbf{d} \right)\text{,}\quad\text{for~all~}\mathbf{d} \in \mathcal{D}\text{.}\]
Intuitively, \(\mathbf{s}^{x}\left( \mathbf{d} \right)\) is the `score'
which alternative \(x\) receives from the profile \(\mathbf{d}\); each
voter who sends the signal \(v\) contributes \(s_{v}^{x}\) `points' to
this score. The alternative with the highest score wins.

For example, plurality vote is a scoring rule with
\(\mathcal{V} = \mathcal{X}\), and \(\mathcal{R} = \mathbb{Z}\), and
\(s_{v}^{x} = 1\) if \(x = v\), while \(s_{v}^{x} = 0\) if \(x \neq v\).
Approval vote is a scoring rule, where \(\mathcal{V}\) is the set of all
subsets of \(\mathcal{X}\), and \(\mathcal{R} = \mathbb{Z}\), and
\(s_{v}^{x} = 1\) if \(x \in v\), while \(s_{v}^{x} = 0\) if
\(x \notin v\). The Borda rule is a scoring rule where \(\mathcal{V}\)
is the set of all strict preference orders over \(\mathcal{X}\), and
\(\mathcal{R} = \mathbb{Z}\), and \(s_{v}^{x} = r\) if \(x\) is ranked
\(r\)th place from the bottom in the preference order \(v\). If
\(\left( \mathcal{Y},d \right)\) is a metric space, and
\(\mathcal{X},\mathcal{V} \subseteq \mathcal{Y}\), then the median rule
is the scoring rule where \(\mathcal{R} = \mathbb{R}\), and
\(s_{v}^{x} = - d(x,v)\) for all \(x \in \mathcal{X}\) and
\(v \in \mathcal{V}\). This rule picks the element(s) of \(\mathcal{X}\)
which minimize the average distance to the signals sent by the voters.
(If \(\mathcal{X} = \mathcal{V} \subseteq \mathbb{R}\) with the standard
metric, then this is the usual notion of the median of a collection of
real numbers.)~For example, the Kemeny rule is a median rule where
\(\mathcal{X} = \mathcal{V}\) is the set of all strict preference orders
over some set \(\mathcal{A}\) of alternatives, and \(d\) is the Kendall
metric on \(\mathcal{X}\) (so \(d(x,v)\) is the number of pairwise
comparisons on which the orderings \(x\) and \(v\) disagree).

There are also several scoring rules where \(\mathcal{R} = \mathbb{R}\)
and \(\mathcal{V}\) is some subset of \(\mathbb{R}^{\mathcal{X}}\), and
\(s_{x}^{\mathbf{v}} ≔ v_{x}\) for all \(\mathbf{v} \in \mathcal{V}\)
and \(x \in \mathcal{X}\). Range voting is obtained by setting
\(\mathcal{V} ≔ \lbrack 0,1\rbrack^{\mathcal{X}}\) (Smith, 2000,
Gaertner and Xu, 2012). Relative utilitarianism is obtained by setting
\(\mathcal{V} ≔ \left\{ \mathbf{v} \in \lbrack 0,1\rbrack^{\mathcal{X}};\min_{x \in \mathcal{X}}v_{x} = 0\text{~and~}\max_{x \in \mathcal{X}}v_{x} = 1 \right\}\)
(Dhillon, 1998, Dhillon and Mertens, 1999). Finally, cumulative voting
is obtained by setting
\(\mathcal{V} ≔ \left\{ \mathbf{v} \in \lbrack 0,1\rbrack^{\mathcal{X}};\sum_{x \in \mathcal{X}}v_{x} = 1 \right\}\)
(Cooper, 2007, Saari, 2010).

Now suppose we first apply one scoring rule \(F_{a}\), and then use a
second scoring rule \(F_{b}\) only to break any ties which arise in
\(F_{a}\). This can be modeled as an \(\mathbb{R}_{lex}^{2}\)-valued
scoring rule. For example, the procedure, `first apply approval vote;
then break any ties using the Borda rule' can be modeled by defining
\(\mathcal{V} = \mathcal{V}_{a} \times \mathcal{V}_{b}\), where
\(\mathcal{V}_{a} ≔ \left\{ \text{all~subsets~of~}\mathcal{X} \right\}\)
and
\(\mathcal{V}_{b} ≔ \left\{ \text{all~preference~orders~over~}\mathcal{X} \right\}\).
For any
\(\left( v_{a},v_{b} \right) \in \mathcal{V}_{a} \times \mathcal{V}_{b}\),
we set
\(s\left( v_{a},v_{b} \right) ≔ \left( s_{a}\left( v_{a} \right),s_{b}\left( v_{b} \right) \right) \in \mathbb{R}_{lex}^{2}\),
where \(s_{a}\) is the approval score system and \(s_{b}\) is the Borda
score system (as described above).

Hahn's Embedding Theorem says that any linearly ordered abelian group is
isomorphic to an ordered subgroup of a lexicographically ordered vector
space \(\mathbb{R}^{\mathcal{I}}\), where \(\mathcal{I}\) is a (possibly
infinite) linearly ordered set (Hausner and Wendel, 1952). Thus, any
scoring rule can be interpreted as a (possibly infinite) chain of
real-valued scoring rules, with each one acting as a tie-breaker for the
prior ones.

Let \(\left( \mathcal{R}, + , > \right)\) be a linearly ordered abelian
group. Let
\(\mathcal{D} \subseteq \mathbb{N}^{\langle\mathcal{V}\rangle}\) be a
domain of profiles. An \(\mathcal{R}\)-valued balance system on
\(\left( \mathcal{X},\mathcal{V},\mathcal{D} \right)\) is an
\(\mathcal{X}^{2}\)-indexed collection
\(\mathsf{B} ≔ \left\{ \mathbf{b}^{x,y} \right\}_{x,y \in \mathcal{X}} \subset \mathcal{R}^{\mathcal{V}}\)
such that \(\mathbf{b}^{x,y} = - \mathbf{b}^{y,x}\) for all
\(x,y \in \mathcal{X}\) (in particular, \(\mathbf{b}^{x,x} = 0\) for all
\(x \in \mathcal{X}\)), and such that,
(1)\[\max\limits_{x \in \mathcal{X}}\min\limits_{y \in \mathcal{X}}\mathbf{b}^{x,y}\left( \mathbf{d} \right) \geq 0\text{,}\quad\text{for~all~}\mathbf{d} \in \mathcal{D}\text{.}\]
We then define the balance
rule\(\left. F_{\mathsf{B}}:\mathcal{D}\rightrightarrows\mathcal{X} \right.\)
as follows: for all \(\mathbf{d} \in \mathcal{D}\) and
\(x \in \mathcal{X}\), we let
\(x \in F_{\mathsf{B}}\left( \mathbf{d} \right)\) if and only if
\(\mathbf{b}^{x,y}\left( \mathbf{d} \right) \geq 0\) for all
\(y \in \mathcal{X}\). (The condition (1) is equivalent to stipulating
that
\(F_{\mathsf{B}}\left( \mathbf{d} \right) \neq \varnothing\)
for all \(\mathbf{d} \in \mathcal{D}\).) Example~1

Let \(\mathsf{S} = \left\{ \mathbf{s}^{x} \right\}_{x \in \mathcal{X}}\)
be an \(\mathcal{R}\)-valued score system on
\(\left( \mathcal{X},\mathcal{V} \right)\). For all
\(x,y \in \mathcal{X}\), define
\(\mathbf{b}^{x,y} ≔ \mathbf{s}^{x} - \mathbf{s}^{y} \in \mathcal{R}^{\mathcal{V}}\),
to obtain a balance system
\(\mathsf{B} ≔ \left\{ \mathbf{b}^{x,y} \right\}_{x,y \in \mathcal{X}}\).
Then
\(F_{\mathsf{B}}\left( \mathbf{n} \right) = F_{\mathsf{S}}\left( \mathbf{n} \right)\)
for all
\(\mathbf{n} \in \mathbb{N}^{\langle\mathcal{V}\rangle}\).\textsuperscript{11}11This
example also appears in Myerson (1995, Section 4, p. 65).\(\quad ♦\)

However, not every balance rule is a scoring rule. Example~2

Let \(\mathcal{V} = \mathcal{X} = \left\{ 1,2,3 \right\}\); thus, there
are three alternatives, and each voter votes for exactly one of these
alternatives. Consider the following rule, defined on
\(\mathcal{D} = \mathbb{N}^{\langle\mathcal{V}\rangle}\). If alternative
1 receives at least half of all votes cast, then alternative 1 wins.
Otherwise, the winner is whichever of alternatives 2 or 3 receives more
votes, with multiple winners in the case of a tie. Formally, given any
profile
\(\mathbf{d} = \left( d_{1},d_{2},d_{3} \right) \in \mathbb{N}^{\langle\mathcal{V}\rangle}\),
if \(d_{1} \geq d_{2} + d_{3}\), then 1 wins. Otherwise, if
\(d_{1} \leq d_{2} + d_{3}\), then 2 wins if \(d_{2} \geq d_{3}\),
whereas \(d_{3}\) wins if \(d_{2} \leq d_{3}\). (If
\(d_{1} = d_{2} + d_{3}\) and/or \(d_{2} = d_{3}\), then there are
multiple winners.) For example, if \(\mathbf{d} = (45,35,30)\), then
\(F\left( \mathbf{d} \right) = \left\{ 2 \right\}\).

One possible rationale for \(F\) is that alternative 1 is a (disfavored)
status quo, while alternatives 2 and 3 are two possible reforms. The
status quo will only be maintained if it receives a clear majority
mandate; otherwise the more popular of the two reforms will be chosen.
It is easy to see that \(F\) is a balance rule, with
\(\mathbf{b}^{1,2} = \mathbf{b}^{1,3} ≔ (1, - 1, - 1)\) and
\(\mathbf{b}^{2,3} ≔ (0,1, - 1)\). But in Appendix~A, we show that \(F\)
is not a scoring rule. \(\quad ♦\)

Let \(\mathcal{D} \subseteq \mathbb{N}^{\langle\mathcal{V}\rangle}\) be
a domain and let
\(\left. F:\mathcal{D}\rightrightarrows\mathcal{X} \right.\) be a voting
rule. We will say that \(\mathcal{D}\) is weakly additive for \(F\) if,
for any \(\mathbf{n},\mathbf{m} \in \mathcal{D}\), if
\(F\left( \mathbf{n} \right) \cap F\left( \mathbf{m} \right) \neq \varnothing\),
then \(\mathbf{n} + \mathbf{m} \in \mathcal{D}\). We will say that the
rule \(\left. F:\mathcal{D}\rightrightarrows\mathcal{X} \right.\)
satisfies reinforcement if \(\mathcal{D}\) is weakly additive for \(F\)
and the following is true: for any
\(\mathbf{n},\mathbf{m} \in \mathcal{D}\), if
\(F\left( \mathbf{n} \right) \cap F\left( \mathbf{m} \right) \neq \varnothing\),
then
\(F\left( \mathbf{n} + \mathbf{m} \right) = F\left( \mathbf{n} \right) \cap F\left( \mathbf{m} \right)\).\textsuperscript{12}12This
is slightly weaker than the usual definition of ``reinforcement'',
because we do not require the domain \(\mathcal{D}\) to be fully closed
under addition---we only require it to be weakly additive. Of course, in
almost all of the ``natural'' examples, \(\mathcal{D}\)~will be closed
under addition. But we allow the possibility that \(\mathcal{D}\) is not
additively closed, in order to accommodate certain interesting examples,
such as the ``Condorcet'' balance rule of Example~3 below. Here, the
profile \(\left( \mathbf{n} + \mathbf{m} \right)\) represents a union of
two disjoint subgroups, represented by profiles \(\mathbf{n}\) and
\(\mathbf{m}\). Reinforcement says: if \(x \in \mathcal{X}\) and both
\(\mathbf{n}\) and \(\mathbf{m}\) endorse \(x\)
(i.e.~\(x \in F\left( \mathbf{n} \right)\) and
\(x \in F\left( \mathbf{m} \right)\)), then we should have
\(x \in F\left( \mathbf{n} + \mathbf{m} \right)\). Furthermore, in this
case, \(F\left( \mathbf{n} + \mathbf{m} \right)\) should consist of only
those \(x \in \mathcal{X}\) which receive this joint endorsement. We now
come to our first main result: Theorem~1

Let \(\mathcal{X}\) and \(\mathcal{V}\) be arbitrary sets, let
\(\mathcal{D} \subseteq \mathbb{N}^{\langle\mathcal{V}\rangle}\) be any
domain, and let
\(\left. F:\mathcal{D}\rightrightarrows\mathcal{X} \right.\) be a
variable population anonymous voting rule for which \(\mathcal{D}\) is
weakly additive. Then \(F\) satisfies reinforcement if and only if \(F\)
is a balance rule.

The proof of Theorem~1 is in Appendix~C, but we will sketch the main
idea here. The first observation is that \(F\) satisfies reinforcement
if and only if, for all \(x \in \mathcal{V}\), the set
\(\mathcal{C}_{x} ≔ \left\{ \mathbf{d} \in \mathcal{D};x \in F\left( \mathbf{d} \right) \right\}\)
is a conoid---i.e.~it is closed under
addition.\textsuperscript{13}13This is closely related to one of the
main insights of Saari (1990), although that paper deals mainly with
weak reinforcement. I thank the referee for pointing out this
connection. It follows that the set
\(\mathcal{P}_{x,y} ≔ \mathcal{C}_{x} - \mathcal{C}_{y}\) is a conoid,
for all \(x,y \in \mathcal{X}\). Such a conoid can be used to define a
homogeneous preorder on \(\mathbb{Z}^{\langle\mathcal{V}\rangle}\). If
\(\mathcal{O}_{x,y} \subset \mathbb{Z}^{\langle\mathcal{V}\rangle}\) is
the set of elements indifferent to \(\mathbf{0}\) in this preorder, then
\(\mathcal{O}_{x,y}\) is itself a subgroup of
\(\mathbb{Z}^{\langle\mathcal{V}\rangle}\), and the quotient group
\(\mathcal{R}_{x,y} ≔ \mathbb{Z}^{\langle\mathcal{V}\rangle}/\mathcal{O}_{x,y}\)
inherits an order which makes it a linearly ordered abelian group. The
quotient map
\(\left. \mathbb{Z}^{\langle\mathcal{V}\rangle}\longrightarrow\mathcal{O}_{x,y} \right.\)
defines an \(\mathcal{R}_{x,y}\)-valued vector
\(\mathbf{b}^{x,y} \in \mathcal{R}_{x,y}^{\mathcal{V}}\). We repeat this
exercise for all \(x,y \in \mathcal{X}\), and then amalgamate all of
\(\left\{ \mathcal{R}_{x,y} \right\}_{x,y \in \mathcal{X}}\) into a
single linearly ordered group \(\mathcal{R}\). The vectors
\(\left\{ \mathbf{b}^{x,y} \right\}_{x,y \in \mathcal{X}}\) then define
an \(\mathcal{R}\)-valued balance system \(\mathsf{B}\), such that
\(F_{\mathsf{B}} = F\).

Example~2 showed that not every balance rule is a scoring rule, even if
we require \(\mathcal{D} = \mathbb{N}^{\langle\mathcal{V}\rangle}\).
Thus, we must add some other hypotheses to reinforcement to characterize
scoring rules. Let \(\Pi_{\mathcal{V}}\) be the group of all
permutations of \(\mathcal{V}\). For any
\(\mathbf{n} \in \mathbb{N}^{\langle\mathcal{V}\rangle}\) and
\(\pi \in \Pi_{\mathcal{V}}\), we define
\(\pi\left( \mathbf{n} \right) ≔ \mathbf{m}\), where
\(m_{v} ≔ n_{\pi^{- 1}{(v)}}\) for all \(v \in \mathcal{V}\). Let
\(\Pi_{\mathcal{X}}\) be the group of all permutations of
\(\mathcal{X}\). A voting rule
\(\left. F:\mathcal{D}\rightrightarrows\mathcal{X} \right.\) is neutral
if there exists a group homomorphism
\(\left. \nu:\Pi_{\mathcal{X}}\longrightarrow\Pi_{\mathcal{V}} \right.\)
(the neutralizer) such that, for all \(\pi \in \Pi_{\mathcal{X}}\), if
\(\widetilde{\pi} ≔ \nu(\pi)\), then the domain \(\mathcal{D}\) is
\(\widetilde{\pi}\)-invariant, and
\(F\left( \widetilde{\pi}\left( \mathbf{d} \right) \right) = \pi\left( F\left( \mathbf{d} \right) \right)\)
for all \(\mathbf{d} \in \mathcal{D}\). Thus, every alternative in
\(\mathcal{X}\) is treated equally: for any \(x,y \in \mathcal{X}\), and
every profile \(\mathbf{d} \in \mathcal{D}\) such that
\(x \in F\left( \mathbf{d} \right)\), there exists some permutation
\(\mathbf{d}'\) of \(\mathbf{d}\) such that
\(y \in F\left( \mathbf{d}' \right)\).

For any \(\pi \in \Pi_{\mathcal{V}}\) and any
\(\mathbf{r} \in \mathcal{R}^{\mathcal{V}}\), we define
\(\mathbf{r}\pi \in \mathcal{R}^{\mathcal{V}}\) by
\(\left( \mathbf{r}\pi \right)_{v} = r_{\pi{(v)}}\) for all
\(v \in \mathcal{V}\). Let
\(\left. \nu:\Pi_{\mathcal{X}}\longrightarrow\Pi_{\mathcal{V}} \right.\)
be a homomorphism. A score system
\(\mathsf{S} = \left\{ \mathbf{s}^{x} \right\}_{x \in \mathcal{X}}\) is
\(\nu\)-neutral if, for all \(\pi \in \Pi_{\mathcal{X}}\) and
\(x,y \in \mathcal{X}\), if \(\pi(y) = x\) and
\(\widetilde{\pi} ≔ \nu(\pi)\), then
\(\mathbf{s}^{x}\,\widetilde{\pi} = \mathbf{s}^{y}\). Except for median
rules, all the scoring rules mentioned above have neutral score systems,
with the obvious neutralizers.

A domain
\(\mathcal{D} \subseteq \mathbb{N}^{\langle\mathcal{V}\rangle}\) is a
cone if \(\mathbf{d}_{1} + \mathbf{d}_{2} \in \mathcal{D}\) whenever
\(\mathbf{d}_{1},\mathbf{d}_{2} \in \mathcal{D}\), and also,
\(\mathbf{d} \in \mathcal{D}\) whenever
\(n\,\mathbf{d} \in \mathcal{D}\) for some \(n \in \mathbb{N}\). For
example, the universal domain \(\mathbb{N}^{\langle\mathcal{V}\rangle}\)
itself is a cone. Also, if
\(\left. f:\mathbb{R}^{\langle\mathcal{V}\rangle}\longrightarrow\mathbb{R} \right.\)
is a linear function, then the sets
\(\left\{ \mathbf{n} \in \mathbb{N}^{\langle\mathcal{V}\rangle};f\left( \mathbf{n} \right) \geq 0 \right\}\)
and
\(\left\{ \mathbf{n} \in \mathbb{N}^{\langle\mathcal{V}\rangle};f\left( \mathbf{n} \right) = 0 \right\}\)
are cones. The intersection of any collection of cones is a cone. Here
is our second main result: Theorem~2

Let \(\mathcal{X}\) be a finite set, let \(\mathcal{V}\) be any set, let
\(\mathcal{D} \subseteq \mathbb{N}^{\langle\mathcal{V}\rangle}\) be a
cone, and let
\(\left. F:\mathcal{D}\rightrightarrows\mathcal{X} \right.\) be a
variable population anonymous voting rule. Then \(F\) is neutral and
satisfies reinforcement if and only if \(F\) is a scoring rule with a
neutral score system.

Thus, combining reinforcement with neutrality yields a scoring rule with
a neutral score system. The next example shows why it is necessary that
\(\mathcal{D}\) be a cone in Theorem~2. Example~3

Let \(\mathcal{V}\) be the set of all nonstrict preference orders over
\(\mathcal{X}\). For all \(x,y \in \mathcal{X}\) and
\(v \in \mathcal{V}\), define \(b_{v}^{x,y} ≔ 1\) if \(v\) prefers \(x\)
to \(y\), while \(b_{v}^{x,y} ≔ - 1\) if \(v\) prefers \(y\) to \(x\),
and \(b_{v}^{x,y} ≔ 0\) if \(v\) is indifferent between \(x\) and \(y\).
Then \(F_{\mathsf{B}}\) is the Condorcet rule: for any
\(\mathbf{n} \in \mathbb{N}^{\langle\mathcal{V}\rangle}\), we have
\(x \in F\left( \mathbf{n} \right)\) if and only if \(x\) is a Condorcet
winner in the profile \(\mathbf{n}\) (i.e.~for any other
\(y \in \mathcal{X}\), at least as many voters strictly prefer \(x\)
over \(y\) as the number who strictly prefer \(y\) over \(x\)).
Unfortunately, \(F\left( \mathbf{n} \right) = \varnothing\)
for some \(\mathbf{n} \in \mathbb{N}^{\langle\mathcal{V}\rangle}\) (the
`Condorcet paradox'). Let
\(\mathcal{D} \subset \mathbb{N}^{\langle\mathcal{V}\rangle}\) be the
set of all profiles having a Condorcet winner. Then
\(\left. F_{\mathsf{B}}:\mathcal{D}\rightrightarrows\mathcal{X} \right.\)
is a neutral balance rule. However, \(\mathcal{D}\) is not a cone (it is
not closed under addition), so Theorem~2 does not apply. And indeed,
\(F_{\mathsf{B}}\) is not a scoring rule.

However, if we restrict \(F_{\mathsf{B}}\) to a smaller domain
\(\mathcal{D}'\) which is a cone, then Theorem~2 implies that
\(F_{\mathsf{B}}\) will be equivalent to a scoring rule on
\(\mathcal{D}'\). For example, let
\(\left. \alpha:\mathcal{X}\longrightarrow\mathbb{R} \right.\) be an
injection. A preference order \(\succ\) on \(\mathcal{X}\)single-peaked
with respect to \(\alpha\) if there is some `ideal point'
\(z_{\succ} \in \mathcal{X}\) such that, for all
\(x,y \in \mathcal{X}\), we have \(y \succ x\) if and only if either
\(\alpha(x) < \alpha(y) < \alpha\left( z_{\succ} \right)\) or
\(\alpha\left( z_{\succ} \right) > \alpha(y) > \alpha(x)\). Let
\(\mathcal{V}_{\alpha} \subset \mathcal{V}\) be the set of
\(\alpha\)-single-peaked preferences, and let
\(\mathcal{D}_{\alpha} ≔ \mathbb{N}^{\langle\mathcal{V}_{\alpha}\rangle}\)
be the set of \(\alpha\)-single-peaked profiles. Then
\(\mathcal{D}_{\alpha}\) is a cone. Furthermore,
\(\mathcal{D}_{\alpha} \subset \mathcal{D}\)---that is, every
single-peaked profile has a Condorcet winner Black (1958).

Since the Condorcet rule is neutral, and \(\mathcal{D}_{\alpha}\) is a
cone, Theorem~2 implies that the Condorcet rule is a scoring rule when
restricted to \(\mathcal{D}_{\alpha}\). Indeed, for all
\(v \in \mathcal{V}\), let \(z_{v} \in \mathcal{X}\) be the `ideal
point' of \(v\), and then, for all \(x \in \mathcal{X}\), define
\(s_{v}^{x} ≔ - \left| \alpha\left( z_{v} \right) - \alpha(x) \right|\).
Thus, the scoring rule \(F_{\mathsf{S}}\) picks the element of
\(\mathcal{X}\) which minimizes the average distance to the ideal points
of the voters---namely, the median of the voter's ideal points. This is
the Condorcet winner Black (1958). Thus, \(F_{\mathsf{S}}\) is
equivalent to the Condorcet rule on \(\mathcal{D}_{\alpha}\).
\(\quad ♦\)

Despite the hypothesis of neutrality in Theorem~2, not all scoring rules
are neutral. For example, abstract median rules are generally not
neutral (but see Example~6 below). Here are two other interesting
examples of non-neutral scoring rules. These rules also illustrate how
non-real score systems can be used to encode structural features
(e.g.~tiebreaking rules, multistage tournaments) which cannot be encoded
with real-valued scoring systems. This shows that non-real score systems
have genuine normative interest, and are not merely technical
curiosities. Example~4

Let \(\mathcal{X} = \left\{ 1,2,\ldots,N \right\}\), and let
\(\mathsf{S}\) be a neutral, real-valued score system which defines a
voting rule
\(\left. F_{\mathsf{S}}:\mathbb{N}^{\langle\mathcal{V}\rangle}\rightrightarrows\mathcal{X} \right.\).
Suppose a committee employs \(F_{\mathsf{S}}\), with ties broken by the
chair. We suppose that the chair has fixed, exogenous preferences, and
that except for tie-breaking, the chair does not vote. Let \(G\) denote
this modified rule. Then \(G\) is obviously not neutral (unless the
chair is indifferent). However, in Appendix~A, we show that \(G\) is
still a scoring rule, with a \(\mathbb{R}_{lex}^{N}\)-valued score
system. \(\quad ♦\)

Example~5

Let \(\mathcal{A} = \left\{ a,b,c,d \right\}\) be a set of four
contestants (e.g.~athletes), and let \(\mathcal{V}\) be the set of all
strict preference orders over \(\mathcal{A}\). The space \(\mathcal{X}\)
of social alternatives will be a subset of \(\mathcal{V}\). To be
precise, let \(\mathcal{X} \subset \mathcal{V}\) consist of all
preference orders on \(\mathcal{A}\) whose top two contestants contain
exactly one element from \(\left\{ a,b \right\}\) and exactly one
element from \(\left\{ c,d \right\}\). (For example,
`\(a \succ c \succ d \succ b\)' is in \(\mathcal{X}\), but
`\(a \succ b \succ c \succ d\)' is not in \(\mathcal{X}\). Neither is
`\(c \succ d \succ a \succ b\)'.) Consider the following `tournament'
rule (used in some sports). Let \(w_{ab}\) be the winner between \(a\)
and \(b\) in a pairwise vote, and let \(\ell_{ab}\) be the loser. Let
\(w_{cd}\) be the winner between \(c\) and \(d\) in a pairwise vote, and
let \(\ell_{cd}\) be the loser. Let \(G\) (`Gold') denote the winner
between \(w_{ab}\) and \(w_{cd}\) in a pairwise vote, and let \(S\)
(`Silver') denote the loser. Let \(B\) (`Bronze') denote the winner
between \(\ell_{ab}\) and \(\ell_{cd}\) in a pairwise vote and let \(L\)
(`Lead') denote the loser. Then the rule \(F\) selects the preference
order `\(G \succ S \succ B \succ L\)' in \(\mathcal{X}\).

If there is a tie in one or more of the four pairwise votes described
above, then \(F\) selects all preference orders in \(\mathcal{X}\) which
are consistent with breaking these ties in an arbitrary way. (For
example, if \(a\) ties \(b\) but all other pairwise votes are non-tied,
then we consider the preference order that would result from the choice
\(w_{ab} = a\) and \(\ell_{ab} = b\), and also the preference order that
would result from the choice \(w_{ab} = b\) and \(\ell_{ab} = a\).)

It is easy to see that \(F\) satisfies reinforcement (because pairwise
votes satisfy reinforcement). But \(F\) is not neutral. For example,
consider the permutation \(\pi \in \Pi_{\mathcal{X}}\) which switches
the bottom two entries in each preference order in \(\mathcal{X}\).
There is no permutation \(\pi' \in \Pi_{\mathcal{V}}\) such that
\(F\left\lbrack \pi'\left( \mathbf{d} \right) \right\rbrack = \pi\left\lbrack F\left( \mathbf{d} \right) \right\rbrack\)
for all \(\mathbf{d} \in \mathbb{N}^{\langle\mathcal{V}\rangle}\).
Nevertheless, in Appendix~A we show that \(F\) is a scoring rule,
defined by a \(\mathbb{R}_{lex}^{4}\)-valued scoring system. \(\quad ♦\)

These examples illustrate that the neutrality condition in Theorem~2 is
sufficient, but not necessary, for a balance rule to be a scoring rule.
Lemma~C.4 (in Appendix~C) provides a simple condition which is both
necessary and sufficient for a balance rule to be a scoring rule.
However, this condition is rather technical and has no obvious normative
interpretation. In Section~3, we will explore the question of neutrality
further.

We finish this section with one final remark. Theorem~2 considers a rule
\(F\) defined on a domain
\(\mathcal{D} \subseteq \mathbb{N}^{\langle\mathcal{V}\rangle}\). Given
a score system \(\mathsf{S}\), we have
\(F_{\mathsf{S}}\left( \mathbf{n} \right) \neq \varnothing\)
for all \(\mathbf{n} \in \mathbb{N}^{\langle\mathcal{V}\rangle}\); thus,
the scoring representation offers a way to `extend' \(F\) from
\(\mathcal{D}\) to all of \(\mathbb{N}^{\langle\mathcal{V}\rangle}\).
However, we might still wish to restrict \(F_{\mathsf{S}}\) to the
smaller domain \(\mathcal{D}\). For example, suppose we have adopted
\(F\) on the basis of certain normative criteria which only make sense
inside \(\mathcal{D}\) (e.g.~Condorcet consistency). If
\(\mathbf{n} \in \mathbb{N}^{\langle\mathcal{V}\rangle} \smallsetminus \mathcal{D}\),
then we might regard the value of
\(F_{\mathsf{S}}\left( \mathbf{n} \right)\) as a meaningless artifact;
the normative criteria which justify \(F\) inside \(\mathcal{D}\) do not
apply at \(\mathbf{n}\).\textsuperscript{14}14A similar remark applies
to a balance rule
\(\left. F_{\mathsf{B}}:\mathcal{D}\rightrightarrows\mathcal{X} \right.\).
There may exist other profiles
\(\mathbf{n} \in \mathbb{N}^{\langle\mathcal{V}\rangle} \smallsetminus \mathcal{D}\)
which satisfy the nontriviality condition (1), so that
\(F_{\mathsf{B}}\left( \mathbf{n} \right) \neq \varnothing\).
But the fact that \(F_{\mathsf{B}}\) is well-defined outside
\(\mathcal{D}\) does not imply that society is normatively compelled to
apply \(F_{\mathsf{B}}\) outside \(\mathcal{D}\).

\protect\phantomsection\label{s000015}
3Generalized neutrality

We have defined a voting rule \(F\) to be neutral if it is equivariant
under all permutations of \(\mathcal{X}\). However, there are many
voting rules which fail to satisfy this stringent condition, but which
are still equivariant under some permutations of
\(\mathcal{X}\)---enough that there is still a sense in which these
rules treat all elements of \(\mathcal{X}\) the same. In this section,
we will develop a more general notion of `neutrality' which encompasses
such rules. The two results of this section then show how the neutrality
properties of a voting rule \(F\) can be `encoded' in a balance system
or scoring system for \(F\). They are also key technical steps in the
proof of Theorem~2.

Let \(\Pi_{\mathcal{X}}' \subseteq \Pi_{\mathcal{X}}\) be any subgroup
of permutations of \(\mathcal{X}\), and let
\(\left. \nu:\Pi_{\mathcal{X}}'\longrightarrow\Pi_{\mathcal{V}} \right.\)
be a group homomorphism (a neutralizer). For any domain
\(\mathcal{D} \subseteq \mathbb{N}^{\langle\mathcal{V}\rangle}\), a
voting rule \(\left. F:\mathcal{D}\rightrightarrows\mathcal{X} \right.\)
is \(\mathsf{\nu}\)-neutral if, for all \(\pi \in \Pi_{\mathcal{X}}'\),
if \(\widetilde{\pi} = \nu(\pi)\), then \(\mathcal{D}\) is
\(\widetilde{\pi}\)-invariant and
\(F\left( \widetilde{\pi}\left( \mathbf{d} \right) \right) = \pi\left( F\left( \mathbf{d} \right) \right)\)
for all \(\mathbf{d} \in \mathcal{D}\).Example~6

Let \(\mathcal{X} = \mathcal{V}\) be a metric space, with metric \(d\).
A self-isometry is a function
\(\left. \pi:\mathcal{X}\longrightarrow\mathcal{X} \right.\) such that
\(d\left\lbrack \pi(x),\pi(y) \right\rbrack = d(x,y)\) for all
\(x,y \in \mathcal{X}\). Let \(\Pi_{\mathcal{X}}'\) be the group of all
self-isometries of \(\mathcal{X}\), and let
\(\left. \nu:\Pi_{\mathcal{X}}'\longrightarrow\Pi_{\mathcal{X}} = \Pi_{\mathcal{V}} \right.\)
be the inclusion map. Then the median rule on \(\mathcal{X}\) is
\(\nu\)-neutral. \(\quad ♦\)

Example~7

The Kemeny rule is not neutral. To see this, let \(\mathcal{A}\) be a
set of candidates, and let the set \(\mathcal{X}\) of social
alternatives be the space of all strict preference orders on
\(\mathcal{A}\). Let \(\mathcal{V} = \mathcal{X}\) (so people vote by
declaring their complete preference order over \(\mathcal{A}\).) If
\(\left| \mathcal{A} \right| = N\), then we can represent
\(\mathcal{X}\) as the space of all bijections
\(\left. x:\mathcal{A}\longrightarrow\lbrack 1\ldots N\rbrack \right.\).
Let \(\Pi_{\mathcal{A}}\) be the group of all permutations of
\(\mathcal{A}\), and let \(\Pi_{N}\) be the group of all permutations of
\(\lbrack 1\ldots N\rbrack\). Then \(\Pi_{\mathcal{A}}\) acts on
\(\mathcal{X}\) by ``composition on the right'', while \(\Pi_{N}\) acts
on \(\mathcal{X}\) by ``composition on the left''. That is, given any
\(\alpha \in \Pi_{\mathcal{A}}\), any \(x \in \mathcal{X}\) and any
\(\beta \in \Pi_{N}\), we can define
\(\alpha^{\ast}(x) ≔ x \circ \alpha\) and
\(\beta^{\ast}(x) ≔ \beta \circ x\), to obtain permutations
\(\alpha^{\ast},\beta^{\ast} \in \Pi_{\mathcal{X}}\).

The Kemeny rule is not equivariant under \(\beta^{\ast}\), for any
nontrivial \(\beta \in \Pi_{N}\). (For example, suppose \(\beta(1) = 2\)
and \(\beta(2) = 1\); then \(\beta^{\ast}\) switches the top two
elements of every preference order in \(\mathcal{X}\). But there is no
permutation of \(\mathcal{V}\) which affects the outcome of the Kemeny
rule in this way.) However, let
\(\Pi_{\mathcal{X}}' ≔ \left\{ \alpha^{\ast};\alpha \in \Pi_{\mathcal{A}} \right\}\);
then \(\Pi_{\mathcal{X}}'\) is a subgroup of \(\Pi_{\mathcal{X}}\).
Since \(\Pi_{\mathcal{V}} = \Pi_{\mathcal{X}}\), we can also regard
\(\Pi_{\mathcal{X}}'\) as a subgroup of \(\Pi_{\mathcal{V}}\); let
\(\left. \nu:\Pi_{\mathcal{X}}'\longrightarrow\Pi_{\mathcal{V}} \right.\)
be the inclusion map. Then the Kemeny rule is \(\nu\)-neutral.
\(\quad ♦\)

Example~8

Let \(\mathcal{V} = \mathcal{X} = \left\{ 1,2,3 \right\}\), as in
Example~2. Consider the following rule: \(F\) chooses the plurality
winner (i.e.~the alternative which receives the most votes), except in
the case of ties, which we resolve as follows. Write `\(x \succ y\)' to
mean that, in the case of a two-way plurality tie between \(x\) and
\(y\), the rule \(F\) will select \(x\) as the winner. Then we use the
tiebreaker relations \(1 \succ 2,2 \succ 3\), and \(3 \succ 1\). (In the
case of a 3-way tie, \(F\) selects the outcome
\(\left\{ 1,2,3 \right\}\).)

To represent this as a scoring rule, let
\(\mathcal{R} = \mathbb{R}_{lex}^{2}\). Define \(l ≔ (1,0)\) and
\(\epsilon ≔ (0,1)\); thus, any
\(\left( r_{1},r_{2} \right) \in \mathcal{R}\) can be written as
\(r_{1}l + r_{2}\epsilon\). (Heuristically, think of \(l\) as the number
1, and \(\epsilon\) is a positive `infinitesimal'.)~Then it is easy to
verify that \(F = F_{\mathsf{S}}\), where
\(\mathbf{s}^{1} = (l,0, - \epsilon),\mathbf{s}^{2} = ( - \epsilon,l,0)\),
and \(\mathbf{s}^{3} ≔ (0, - \epsilon,l)\).

This rule is clearly not neutral. For example, let
\(\pi \in \Pi_{\mathcal{X}}\) be the permutation which swaps 1 and 2 but
fixes 3. There is no permutation \(\pi' \in \Pi_{\mathcal{V}}\) such
that
\(F\left\lbrack \pi'\left( \mathbf{d} \right) \right\rbrack = \pi\left\lbrack F\left( \mathbf{d} \right) \right\rbrack\)
for all \(\mathbf{d} \in \mathbb{N}^{\langle\mathcal{V}\rangle}\). (To
see this, let \(\mathbf{d} ≔ (1,1,0)\).)

However, let \(\Pi_{\mathcal{X}}'\) be the subset of
\(\Pi_{\mathcal{X}}\) generated by the permutation
\(\left( 1\rightarrow 2\rightarrow 3\rightarrow 1 \right)\). Since
\(\Pi_{\mathcal{V}} = \Pi_{\mathcal{X}}\), we can also regard
\(\Pi_{\mathcal{X}}'\) as a subgroup of \(\Pi_{\mathcal{V}}\); let
\(\left. \nu:\Pi_{\mathcal{X}}'\longrightarrow\Pi_{\mathcal{V}} \right.\)
be the inclusion map. Then \(F\) is \(\nu\)-neutral. \(\quad ♦\)

The next result says that it is possible to represent a scoring rule
\(F\) with a score system \(\mathsf{S}\) such that \(F\)'s neutrality
properties can be read off explicitly from \(\mathsf{S}\). Proposition~1

Let \(\mathcal{X}\) be a finite set, let \(\mathcal{V}\) be an arbitrary
set, let
\(\mathcal{D} \subseteq \mathbb{N}^{\langle\mathcal{V}\rangle}\) be a
domain, and let
\(\left. F:\mathcal{D}\rightrightarrows\mathcal{X} \right.\) be a
scoring rule. Let \(\Pi_{\mathcal{X}}'\) be any group of permutations of
\(\mathcal{X}\), and let
\(\left. \nu:\Pi_{\mathcal{X}}'\longrightarrow\Pi_{\mathcal{V}} \right.\)
be a homomorphism. Then \(F\) is \(\nu\)-neutral if and only if
\(F = F_{\mathsf{S}}\) for some \(\nu\)-neutral score system
\(\mathsf{S}\).

The corresponding result for balance rules is a bit more complicated. A
balance system \(\mathsf{B}\) is \(\nu\)-neutral if, for all
\(x,y,x',y' \in \mathcal{X}\) and \(\pi \in \Pi_{\mathcal{X}}'\), if
\(x' ≔ \pi^{- 1}(x)\) and \(y' ≔ \pi^{- 1}(y)\) and
\(\widetilde{\pi} = \nu(\pi)\), then
\(\mathbf{b}^{x,y}\widetilde{\pi} = \mathbf{b}^{x',y'}\). A score system
\(\mathsf{S} = \left\{ \mathbf{s}^{x} \right\}_{x \in \mathcal{X}}\) is
\(\nu\)-neutral if, for all \(\pi \in \Pi_{\mathcal{X}}'\) and
\(x,y \in \mathcal{X}\), if \(\pi(y) = x\) and
\(\widetilde{\pi} ≔ \nu(\pi)\), then
\(\mathbf{s}^{x}\,\widetilde{\pi} = \mathbf{s}^{y}\). A voting rule,
balance system, or score system is \(\Pi_{\mathcal{X}}'\)-neutral if it
is \(\nu\)-neutral for some group homomorphism
\(\left. \nu:\Pi_{\mathcal{X}}'\longrightarrow\Pi_{\mathcal{V}} \right.\).
Example~9

Let \(\mathcal{V}\) be the set of preference orders over
\(\mathcal{X}\). Any permutation of \(\mathcal{X}\) acts on
\(\mathcal{V}\) in the obvious manner; let
\(\left. \nu:\Pi_{\mathcal{X}}\longrightarrow\Pi_{\mathcal{V}} \right.\)
be the corresponding homomorphism. Then the Condorcet balance system
(Example~3) is \(\nu\)-neutral. \(\quad ♦\)

Let \(\Pi_{\mathcal{X}}'\) be a group of permutations of
\(\mathcal{X}\). We say that \(\Pi_{\mathcal{X}}'\) is doubly transitive
if, for all \(x,y,x',y' \in \mathcal{X}\) with \(x \neq y\) and
\(x' \neq y'\), there exists some \(\pi \in \Pi_{\mathcal{X}}'\) such
that \(\pi(x) = x'\) and \(\pi(y) = y'\). (For example, the group
\(\Pi_{\mathcal{X}}\) of all permutations of \(\mathcal{X}\) is doubly
transitive.)

A balance system \(\mathsf{B}\) is perfect on the domain \(\mathcal{D}\)
if, for any \(\mathbf{d} \in \mathcal{D}\), any
\(x \in F_{\mathsf{B}}\left( \mathbf{d} \right)\) and any
\(y \in \mathcal{X} \smallsetminus F_{\mathsf{B}}\left( \mathbf{d} \right)\),
we have \(\mathbf{b}^{x,y}\left( \mathbf{d} \right) > 0\). (For
instance: Example~1 is perfect, but Example~2, Example~3 are not
perfect.) Perfection simplifies the computation of
\(F\left( \mathbf{d} \right)\): once we find one point
\(x \in F_{\mathsf{B}}\left( \mathbf{d} \right)\), we have
\(F_{\mathsf{B}}\left( \mathbf{d} \right) = \left\{ y \in \mathcal{X};\mathbf{b}^{x,y}\left( \mathbf{d} \right) = 0 \right\}\).

The next result says that, assuming double-transitivity, it is possible
to represent a balance rule in a manner such that its neutrality
properties can be read off explicitly from its balance system.
Furthermore, we can ensure that this balance system is perfect.
Proposition~2

Let \(\mathcal{X}\) be a finite set, let \(\mathcal{V}\) be an arbitrary
set, let
\(\mathcal{D} \subseteq \mathbb{N}^{\langle\mathcal{V}\rangle}\) be a
domain, and let
\(\left. F:\mathcal{D}\rightrightarrows\mathcal{X} \right.\) be a
balance rule for which \(\mathcal{D}\) is weakly additive. Let
\(\Pi_{\mathcal{X}}'\) be a doubly transitive group of permutations of
\(\mathcal{X}\), and let
\(\left. \nu:\Pi_{\mathcal{X}}'\longrightarrow\Pi_{\mathcal{V}} \right.\)
be a homomorphism. Then \(F\) is \(\nu\)-neutral if and only if
\(F = F_{\mathsf{B}}\) for some \(\nu\)-neutral perfect balance system
\(\mathsf{B}\).

\protect\phantomsection\label{s000020}
4Conclusion

In the setting of anonymous, variable population voting rules, I have
characterized the class of balance rules by a single axiom:
reinforcement. I have also shown that any neutral voting rule which
satisfies reinforcement is a scoring rule. I allowed the underlying
balance system or scoring system to range over an arbitrary linearly
ordered abelian group (rather than the real numbers). This is not just a
technicality; several examples illustrate how this generalization
provides the rules with additional flexibility. Pivato (2012b) discusses
the extent to which real-valued balance or scoring representations of a
voting rule are unique up to scalar multiplication; this has
implications for the ways in which scoring rules can be interpreted as
statistical estimators (Pivato, 2013). Meanwhile, Pivato (2012a) uses
some of the tools developed in this paper to characterize two scoring
rules: range voting and formal utilitarianism.

I will end with some open problems. The conditions of reinforcement and
neutrality are sufficient to obtain a scoring rule (Theorem~2), but
neutrality is not necessary. Are there normatively compelling conditions
which are both necessary and sufficient for a scoring rule? Aside from
Example~2, are there normatively compelling balance rules which are not
scoring rules? Finally, most of the results in this paper assume
\(\mathcal{X}\) is finite---can this hypothesis be eliminated?

Acknowledgments

This paper was written while visiting the Department of Economics at the
Université de Montréal, and the Centre for Philosophy of Natural and
Social Sciences at the London School of Economics. I would like to thank
the UdM, CIREQ, and LSE-CPNSS for their hospitality. I am also grateful
to Matias Nuñez Rodriguez for his comments on a draft of the paper.
Finally, I would like to thank the anonymous referee for many
suggestions which improved the expositional clarity of the paper. This
research was supported by NSERC grant \#262620-2008.

\protect\phantomsection\label{s000025}
Appendix AProofs of some examples

Proof of Example~2

\textsuperscript{15}15I thank the referee for suggesting this simplified
version of the proof. By contradiction, suppose \(F\) was a scoring
rule, with score system
\(\mathsf{S} = \left\{ \mathbf{s}^{1},\mathbf{s}^{2},\mathbf{s}^{3} \right\}\),
where
\(\mathbf{s}^{x} = \left( s_{1}^{x},s_{2}^{x},s_{3}^{x} \right) \in \mathcal{R}^{3}\)
for all \(x \in \mathcal{X} = \left\{ 1,2,3 \right\}\).

Now, \(F(1,0,0) = \left\{ 1 \right\}\); thus, if
\(\mathbf{d} = (1,0,0)\), then we must have
\(\mathbf{s}^{1}\left( \mathbf{d} \right) > \mathbf{s}^{2}\left( \mathbf{d} \right)\),
which means \(s_{1}^{1} > s_{1}^{2}\), which means
\(s_{1}^{2} = s_{1}^{1} - p_{1}\) for some \(p_{1} \in \mathcal{R}\)
with \(p_{1} > 0\). Likewise, \(F(0,1,0) = \left\{ 2 \right\}\); thus,
\(s_{2}^{2} > s_{2}^{1}\), which means \(s_{2}^{1} = s_{2}^{2} - p_{2}\)
for some \(p_{2} > 0\). Finally, \(F(1,1,0) = \left\{ 1,2 \right\}\).
Thus, we must have \(s_{1}^{1} + s_{2}^{1} = s_{1}^{2} + s_{2}^{2}\); in
other words,
\(s_{1}^{1} + s_{2}^{2} - p_{2} = s_{1}^{1} - p_{1} + s_{2}^{2}\), from
which it follows that \(p_{1} = p_{2}\). We conclude that there is a
single value \(p > 0\) such that \(s_{1}^{2} = s_{1}^{1} - p\) and
\(s_{2}^{1} = s_{2}^{2} - p\).

By identical arguments, we can show that there exists \(q > 0\) such
that \(s_{1}^{3} = s_{1}^{1} - q\) and \(s_{3}^{1} = s_{3}^{3} - q\),
and there exists \(r > 0\) such that \(s_{3}^{2} = s_{3}^{3} - r\) and
\(s_{2}^{3} = s_{2}^{2} - r\).

Finally, observe that \(F(2,1,1) = \left\{ 1,2,3 \right\}\). Thus, if
\(\mathbf{d} = (2,1,1)\), then we must have
\(\mathbf{s}^{1}\left( \mathbf{d} \right) = \mathbf{s}^{2}\left( \mathbf{d} \right) = \mathbf{s}^{3}\left( \mathbf{d} \right)\),
which means
\[2s_{1}^{1} + s_{2}^{1} + s_{3}^{1} = 2s_{1}^{2} + s_{2}^{2} + s_{3}^{2} = 2s_{1}^{3} + s_{2}^{3} + s_{3}^{3}\text{.}\]
Substituting in the identities discovered above, this becomes:
\[2s_{1}^{1} + s_{2}^{2} - p + s_{3}^{3} - q = 2s_{1}^{1} - 2p + s_{2}^{2} + s_{3}^{3} - r = 2s_{1}^{1} - 2q + s_{2}^{2} - r + s_{3}^{3}\]
which reduces to \(- p - q = - 2p - r = - 2q - r\). This implies that
\(p = q\), which then implies that \(r = 0\). But \(r > 0\), a
contradiction. □~Example~2

Proof of Example~4

Let \(\mathcal{R} ≔ \mathbb{R}_{lex}^{N}\), and let \(\gg\) denote the
lexicographical order on \(\mathbb{R}_{lex}^{N}\). Recall that
\(\mathcal{X} = \left\{ 1,2,\ldots,N \right\}\). Suppose the committee
chair has the preference order \(1 \prec 2 \prec \cdots \prec N\).
Define the \(\mathbb{R}_{lex}^{N}\)-valued scoring rule
\(\left\{ {\widetilde{\mathbf{s}}}^{x} \right\}_{x \in \mathcal{X}}\) as
follows: for all \(v \in \mathcal{V}\),
\[{\widetilde{s}}_{v}^{1} = \left( s_{v}^{1},0,0,0,\ldots,0 \right) \in \mathbb{R}_{lex}^{N}\text{,}\]\[{\widetilde{s}}_{v}^{2} = \left( s_{v}^{2},1,0,0,\ldots,0 \right) \in \mathbb{R}_{lex}^{N}\text{,}\]\[{\widetilde{s}}_{v}^{3} = \left( s_{v}^{3},1,1,0,\ldots,0 \right) \in \mathbb{R}_{lex}^{N}\text{,}\]\[\vdots\]\[\text{and}\quad{\widetilde{s}}_{v}^{N} = \left( s_{v}^{N},1,1,1,\ldots,1 \right) \in \mathbb{R}_{lex}^{N}\text{.}\]
Then \(F_{\widetilde{\mathsf{S}}} = G\). To see this, let
\(\mathbf{d} \in \mathcal{D}\) be a profile, and let
\(x,y \in \mathcal{X}\). Then
\({\widetilde{\mathbf{s}}}^{x}\left( \mathbf{d} \right)\) and
\({\widetilde{\mathbf{s}}}^{y}\left( \mathbf{d} \right)\) are elements
of \(\mathbb{R}_{lex}^{N}\) whose first (i.e.~lexicographically
dominant) coordinates are \(\mathbf{s}^{x}\left( \mathbf{d} \right)\)
and \(\mathbf{s}^{y}\left( \mathbf{d} \right)\), respectively.
Meanwhile, the next \(x - 1\) coordinates of
\(\mathbf{s}^{x}\left( \mathbf{d} \right)\) are positive, whereas the
next \(y - 1\) coordinates of
\(\mathbf{s}^{y}\left( \mathbf{d} \right)\) are positive, and the
remaining coordinates of both are zero. Thus, if
\(\mathbf{s}^{x}\left( \mathbf{d} \right) > \mathbf{s}^{y}\left( \mathbf{d} \right)\),
then
\({\widetilde{\mathbf{s}}}^{x}\left( \mathbf{d} \right) \gg {\widetilde{\mathbf{s}}}^{y}\left( \mathbf{d} \right)\).
But if
\(\mathbf{s}^{x}\left( \mathbf{d} \right) = \mathbf{s}^{y}\left( \mathbf{d} \right)\),
then
\({\widetilde{\mathbf{s}}}^{x}\left( \mathbf{d} \right) \gg {\widetilde{\mathbf{s}}}^{y}\left( \mathbf{d} \right)\)
if and only if \(x > y\). It follows: if
\(F_{\mathsf{S}}\left( \mathbf{d} \right) = \left\{ x \right\}\), then
\(F_{\widetilde{\mathsf{S}}}\left( \mathbf{d} \right) = \left\{ x \right\}\).
But if
\(\left| F_{\mathsf{S}}\left( \mathbf{d} \right) \right| \geq 2\), then
\(F_{\widetilde{\mathsf{S}}}\left( \mathbf{d} \right)\) is the
numerically largest element of
\(F_{\mathsf{S}}\left( \mathbf{d} \right)\)---that is, the chair's
preferred choice. Hence \(F_{\widetilde{\mathsf{S}}} = G\). □~Example~4

Proof of Example~5

Let \(\mathcal{R} ≔ \mathbb{R}_{lex}^{4}\). For any
\(x \in \mathcal{X}\) and \(v \in \mathcal{V}\), define
\(s_{v}^{x} = \left( {}_{{}_{1}}s_{v}^{x},\,{}_{{}_{2}}s_{v}^{x},{}_{{}_{3}}s_{v}^{x},\,{}_{{}_{4}}s_{v}^{x} \right) \in \mathbb{R}_{lex}^{4}\)
as follows: First, \({}_{{}_{1}}s_{v}^{x} = 1\) if \(v\) agrees with
\(x\)'s ranking of \(a\) vs. \(b\); otherwise
\({}_{{}_{1}}s_{v}^{x} = - 1\). Next, \({}_{{}_{2}}s_{v}^{x} = 1\) if
\(v\) agrees with \(x\)'s ranking of \(c\) vs. \(d\); otherwise
\({}_{{}_{2}}s_{v}^{x} = - 1\). Now suppose the top two choices of \(x\)
are \(g\) and \(h\). Then \({}_{{}_{3}}s_{v}^{x} = 1\) if \(v\) agrees
with \(x\)'s ranking of \(g\) vs. \(h\); otherwise
\({}_{{}_{3}}s_{v}^{x} = - 1\). Finally, suppose the bottom two choices
of \(x\) are \(i\) and \(j\). Then \({}_{{}_{4}}s_{v}^{x} = 1\) if \(v\)
agrees with \(x\)'s ranking of \(i\) vs. \(j\); otherwise
\({}_{{}_{4}}s_{v}^{x} = - 1\).

We claim that \(F_{\mathsf{S}} = F\). To see this, let
\(\mathbf{d} \in \mathbb{N}^{\langle\mathcal{V}\rangle}\) be some
profile. If (for example), a strict majority favors \(a\) over \(b\),
then for all \(x \in \mathcal{X}\), we will have
\({}_{{}_{1}}\mathbf{s}^{x}\left( \mathbf{d} \right) > 0\) if and only
if \(x \in \mathcal{X}_{ab}\), where
\(\mathcal{X}_{ab} ≔ \left\{ x \in \mathcal{X};x\text{~favors~}a\text{~over~}b \right\}\).
Thus,
\(F_{\mathsf{S}}\left( \mathbf{d} \right) \subseteq \mathcal{X}_{ab}\)
(by the lexicographical ordering of \(\mathbb{R}_{lex}^{4}\)).
Furthermore, \({}_{{}_{1}}\mathbf{s}^{x}\left( \mathbf{d} \right)\) has
the same value for all \(x \in \mathcal{X}_{ab}\); thus, the choice
amongst these elements must be made by looking at
\({}_{{}_{2}}\mathbf{s}^{x}\left( \mathbf{d} \right)\). Next, if (for
example) a strict majority favors \(c\) over \(d\), then for all
\(x \in \mathcal{X}_{ab}\), we will have
\({}_{{}_{2}}\mathbf{s}^{x}\left( \mathbf{d} \right) > 0\) if and only
if \(x \in \mathcal{X}_{cd}\), where
\(\mathcal{X}_{cd} ≔ \left\{ x \in \mathcal{X};x\text{~favors~}c\text{~over~}d \right\}\).
Thus,
\(F_{\mathsf{S}}\left( \mathbf{d} \right) \subseteq \mathcal{X}_{ab} \cap \mathcal{X}_{cd}\)
(by the lexicographical ordering of \(\mathbb{R}_{lex}^{4}\)).
Furthermore, \({}_{{}_{2}}\mathbf{s}^{x}\left( \mathbf{d} \right)\) has
the same value for all \(x \in \mathcal{X}_{ab} \cap \mathcal{X}_{cd}\);
thus, the choice amongst these elements must be made by looking at
\({}_{{}_{3}}\mathbf{s}^{x}\left( \mathbf{d} \right)\). Next, if (for
example) a strict majority favors \(c\) over \(a\), then for all
\(x \in \mathcal{X}_{ab} \cap \mathcal{X}_{cd}\), we will have
\({}_{{}_{3}}\mathbf{s}^{x}\left( \mathbf{d} \right) > 0\) if and only
if
\(x \in \mathcal{X}_{ca} ≔ \left\{ x \in \mathcal{X}; x\text{~favors~}c\text{~over~}a \right\}\).
Thus,
\(F_{\mathsf{S}}\left( \mathbf{d} \right) \subseteq \mathcal{X}_{ab} \cap \mathcal{X}_{cd} \cap \mathcal{X}_{ca}\).
Furthermore, \({}_{{}_{3}}\mathbf{s}^{x}\left( \mathbf{d} \right)\) has
the same value for all
\(x \in \mathcal{X}_{ab} \cap \mathcal{X}_{cd} \cap \mathcal{X}_{ca}\);
thus, the choice amongst these elements must be made by looking at
\({}_{{}_{4}}\mathbf{s}^{x}\left( \mathbf{d} \right)\). Finally, if (for
example) a strict majority favors \(b\) over \(d\), then for all
\(x \in \mathcal{X}_{ab} \cap \mathcal{X}_{cd} \cap \mathcal{X}_{ca}\),
we will have \({}_{{}_{4}}\mathbf{s}^{x}\left( \mathbf{d} \right) > 0\)
if and only if
\(x \in \mathcal{X}_{bd} ≔ \left\{ x \in \mathcal{X};x\text{~favors~}b\text{~over~}d \right\}\).
Thus,
\(F_{\mathsf{S}}\left( \mathbf{d} \right) \subseteq \mathcal{X}_{ab} \cap \mathcal{X}_{cd} \cap \mathcal{X}_{ca} \cap \mathcal{X}_{bd} = \left\{ c \succ a \succ b \succ d \right\}\).
□~Example~5

\protect\phantomsection\label{s000030}
Appendix BHomogeneous orders on abelian groups

Let \(\left( \mathcal{R}, + \right)\) be an abelian group, and let
\(( \succeq )\) be a binary relation on \(\mathcal{R}\). We say that
\(( \succeq )\) is homogeneous if, for all \(r,s \in \mathcal{R}\), we
have
\(\left. (r \succeq s)\;\Longleftrightarrow\;(r - s \succeq 0) \right.\).
The positive conoid\textsuperscript{16}16This object is often called the
positive cone. However, we are reserving the term cone for a set which
is divisible. Hence the ungainly term `conoid'. of \(( \succeq )\) is
the set
\(\mathcal{P}_{\succeq} ≔ \left\{ r \in \mathcal{R};r \succeq 0 \right\}\).
The conoid \(\mathcal{P}_{\succeq}\) completely encodes the relation
\(( \succeq )\): for any \(r,s \in \mathcal{R}\), we have
(B.1)\[\left. (r \succeq s)\;\Longleftrightarrow\;\left( r - s \in \mathcal{P}_{\succeq} \right)\text{.} \right.\]
Conversely, given any subset \(\mathcal{P} \subseteq \mathcal{R}\), we
can use formula (B.1) to define a unique homogeneous binary relation
\(( \succeq )\) such that \(\mathcal{P}_{\succeq} = \mathcal{P}\). Thus,
there is a bijective correspondence between homogeneous binary relations
on \(\mathcal{R}\) and subsets of \(\mathcal{R}\).

Let \(\mathcal{P} \subseteq \mathcal{R}\). Recall that \(\mathcal{P}\)
is additively closed if \(p_{1} + p_{2} \in \mathcal{P}\) whenever
\(p_{1},p_{2} \in \mathcal{P}\). We say that \(\mathcal{P}\) is
divisible if, for all \(r \in \mathcal{R}\) and \(n \in \mathbb{N}\), if
\(n\, r \in \mathcal{P}\), then \(r \in \mathcal{P}\) also. Thus,
\(\mathcal{P}\) is a cone if it is both additively closed and divisible.
A preorder conoid is an additively closed subset of \(\mathcal{R}\)
which contains 0. A partial order conoid is an additively closed subset
of \(\mathcal{R}\) which does not contain 0. (Equivalently: for all
\(r \in \mathcal{R}\), at most one of \(r\) or \(- r\) is in
\(\mathcal{P}\)).~A preorder or partial order conoid \(\mathcal{P}\) is
complete if, for all
\(r \in \mathcal{R} \smallsetminus \left\{ 0 \right\}\), at least one of
\(r\) or \(- r\) is in \(\mathcal{P}\). (Any complete preorder or
partial order conoid is divisible.)~A linear order cone is a complete
partial order conoid. (Thus, for all
\(r \in \mathcal{R} \smallsetminus \left\{ 0 \right\}\), exactly one of
\(r \in \mathcal{P}\) or \(- r \in \mathcal{P}\).)~Properties of
\(( \succeq )\) correspond to properties of the conoid
\(\mathcal{P}_{\succeq}\) as follows: Lemma~B.1

Let \(\left( \mathcal{R}, + \right)\) be an abelian group. Let
\(( \succeq )\) be a homogeneous binary relation on \(\mathcal{R}\).(a)

\(( \succeq )\) is transitive iff \(\mathcal{P}_{\succeq}\) is
additively closed. Likewise, \(( \succeq )\) is a preorder iff
\(\mathcal{P}_{\succeq}\) is a preorder conoid. Also, \(( \succeq )\) is
a partial order iff \(\mathcal{P}_{\succeq}\) is a partial order conoid.
Finally, \(( \succeq )\) is a linear order iff \(\mathcal{P}_{\succeq}\)
is a linear order cone.

(b)

Suppose \(( \succeq )\) is a homogeneous preorder; let \(( \succ )\) be
its asymmetric part, and let \(( \approx )\) be its symmetric part. The
set \(\mathcal{O} ≔ \left\{ r \in \mathcal{R};r \approx 0 \right\}\) is
a subgroup of \(\mathcal{R}\) . Let
\(\mathcal{R}' ≔ \mathcal{R}/\mathcal{O}\) be the quotient group, let
\(\left. \phi:\mathcal{R}\longrightarrow\mathcal{R}' \right.\) be the
quotient map, and define
\(\mathcal{P}' ≔ \phi\left( \mathcal{P} \right) \smallsetminus \left\{ 0' \right\}\),
where \(0'\) is the identity element in \(\mathcal{R}'\) . Then
\(\mathcal{P}'\) is a partial order conoid on \(\mathcal{R}'\).

Let \(\left( \succ ' \right)\) be the homogeneous partial order on
\(\mathcal{R}'\) defined by the conoid \(\mathcal{P}'\) . Then for all
\(r,s \in \mathcal{R}\), we have \(\phi(r) \succ '\phi(s)\) if and only
if \(r \succ s\).

Proof

(a)

\(( \succeq )\) is transitive \(\left. \;\Longleftrightarrow\; \right.\)
(for all \(r,s,t \in \mathcal{R}\), if \(r \succeq s\) and
\(s \succeq t\), then \(r \succeq t\))
\(\left. \;\Longleftrightarrow\; \right.\) (for all
\(r,s,t \in \mathcal{R}\), if \((r - s) \in \mathcal{P}_{\succeq}\) and
\((s - t) \in \mathcal{P}_{\succeq}\), then
\((r - t) \in \mathcal{P}_{\succeq}\))
\(\left. \;\Longleftrightarrow\; \right.\) (for all
\(a,b \in \mathcal{R}\), if \(a \in \mathcal{P}_{\succeq}\) and
\(b \in \mathcal{P}_{\succeq}\), then \(a + b \in \mathcal{P}\))
\(\left. \;\Longleftrightarrow\;\mathcal{P}_{\succeq} \right.\) is
additively closed. The proofs of the other statements are similar.

(b)

Observe that \(\mathcal{O} = - \mathcal{P} \cap \mathcal{P}\); thus,
\(\mathcal{O}\) is closed under addition because \(\mathcal{P}\) is
closed under addition. Furthermore, \(\mathcal{O}\) is closed under
negation by definition; thus it is a subgroup of \(\mathcal{R}\).

By definition, \(0' = \phi\left( \mathcal{O} \right)\). The set
\(\mathcal{P}'\) excludes \(0'\) by definition. To see that
\(\mathcal{P}'\) is closed under addition, let
\(p',q' \in \mathcal{P}\); then there exist \(p,q \in \mathcal{P}\) such
that \(\phi(p) = p'\) and \(\phi(q) = q'\). Thus,
\(p + q \in \mathcal{P}\) (because \(\mathcal{P}\) is additively
closed), and \(\phi(p + q) = p' + q'\) (because \(\phi\) is a
homomorphism). To show that \(p' + q' \in \mathcal{P}\), we must show
that \(p' + q' \neq 0'\).

By contradiction, suppose \(p' + q' = 0'\); then
\(p + q \in \mathcal{O}\). Then \(- p - q \in \mathcal{O}\) (because
\(\mathcal{O}\) is closed under negation). Thus,
\(- p - q \in \mathcal{P}\) (because
\(\mathcal{O} \subseteq \mathcal{P}\)). Thus,
\(- q = p + ( - p - q) \in \mathcal{P}\), and
\(- p = q + ( - p - q) \in \mathcal{P}\) (because \(\mathcal{P}\) is
closed under addition). It follows that \(p,q \in \mathcal{O}\). But
then \(p' = 0' = q'\), so \(p',q' \notin \mathcal{P}'\), a
contradiction.

It follows that \(\mathcal{P}'\) is a partial order cone. Finally, for
any \(r,s \in \mathcal{R}\), we have
\(\left. (r \succ s)\;\Longleftrightarrow\;\left( (r - s) \in \mathcal{P} \smallsetminus \mathcal{O} \right)\;\Longleftrightarrow\;\left( \phi(r) - \phi(s) \in \mathcal{P}' \right)\;\Longleftrightarrow\;\left( \phi(r) \succ '\phi(s) \right) \right.\).
□~Lemma~B.1

Next, we extend a classic result about partial orders to the setting of
homogeneous orders. If \(( \succeq )\) and \(\left( \succeq ' \right)\)
are two relations on \(\mathcal{R}\), we say that
\(\left( \succeq ' \right)\)extends\(( \succeq )\) if, for all
\(r,s \in \mathcal{R}\), we have
\(\left. (r \succeq s)\Longrightarrow\left( r \succeq 's \right) \right.\).
An abelian group \(\left( \mathcal{R}, + \right)\) is torsion-free if
\(n\, r \neq 0\) for any
\(n \in \mathbb{Z} \smallsetminus \left\{ 0 \right\}\) and
\(r \in \mathcal{R} \smallsetminus \left\{ 0 \right\}\). Homogeneous
Szpilrajn Lemma

Let \(\left( \mathcal{R}, + \right)\) be a torsion-free abelian group.
Any homogeneous partial order on \(\mathcal{R}\) is extended by a
homogeneous linear order.

Proof

See Fuchs (1950) or Everett (1950).□

Corollary~B.1

Let \(\left( \mathcal{A}, + \right)\) be an abelian group.(a)

If \(\left( \mathcal{A}, + \right)\) is torsion-free, then
\(\mathcal{A}\) has a homogeneous linear ordering.

(b)

Let \(\mathcal{B} \subset \mathcal{A}\) be a divisible proper subgroup.
Then there exists a complete, homogeneous preorder \(( \succeq )\) on
\(\mathcal{A}\) such that
\(\mathcal{B} = \left\{ a \in \mathcal{A};a \approx 0 \right\}\).

Proof

(a)

Let \(a \in \mathcal{A}\) be arbitrary. Let
\(\mathcal{P} ≔ \left\{ n\, a;n \in \mathbb{N} \right\}\). Then
\(\mathcal{P}\) is additively closed and \(0 \notin \mathcal{P}\)
(because \(\mathcal{A}\) is torsion-free), so \(\mathcal{P}\) is a
partial order conoid on \(\mathcal{A}\). If \(( \succ )\) is the
homogeneous partial order defined by \(\mathcal{P}\), then the
Homogeneous Szpilrajn Lemma extends \(( \succ )\) to a homogeneous
linear order on \(\mathcal{A}\).

(b)

Let \(\mathcal{Q} ≔ \mathcal{A}/\mathcal{B}\), and let
\(\left. q:\mathcal{A}\longrightarrow\mathcal{Q} \right.\) be the
quotient map. Then \(\mathcal{Q}\) is a torsion-free group, because
\(\mathcal{B}\) is divisible. Thus, part (a) yields a linear order
\(( > )\) on \(\mathcal{Q}\). Now, define \(( \succeq )\) on
\(\mathcal{A}\) by stipulating that \(a_{1} \succeq a_{2}\) if and only
if \(q\left( a_{1} \right) \geq q\left( a_{2} \right)\). It is easy to
check that \(( \succeq )\) is a complete, homogeneous preorder, and
clearly,
\(\mathcal{B} = \left\{ a \in \mathcal{A};q(a) = 0 \right\} = \left\{ a \in \mathcal{A};a \approx 0 \right\}\).□~Corollary~B.1

\protect\phantomsection\label{s000035}
Appendix CProofs of the main results

Proof of Theorem~1

``\(\Longleftarrow\)'' Suppose \(F\) is a balance rule determined by
some balance system
\(\mathsf{B} = \left\{ \mathbf{b}^{x,y} \right\}_{x,y \in \mathcal{X}}\).
We must show that \(F\) satisfies reinforcement. Let
\(\mathbf{n},\mathbf{m} \in \mathcal{D}\), and suppose
\(F\left( \mathbf{n} \right) \cap F\left( \mathbf{m} \right) \neq \varnothing\).
Then \(\mathbf{n} + \mathbf{m} \in \mathcal{D}\), because
\(\mathcal{D}\) is weakly additive. Now, let
\(x \in F\left( \mathbf{n} \right) \cap F\left( \mathbf{m} \right)\).
Then \(x \in F\left( \mathbf{n} \right)\) and
\(x \in F\left( \mathbf{m} \right)\). Thus, for all
\(y \in \mathcal{X}\), we have
\(\mathbf{s}^{x}\left( \mathbf{n} \right) \geq \mathbf{s}^{y}\left( \mathbf{n} \right)\)
and
\(\mathbf{s}^{x}\left( \mathbf{m} \right) \geq \mathbf{s}^{y}\left( \mathbf{m} \right)\).
Thus,
\[\mathbf{s}^{x}\left( \mathbf{n} + \mathbf{m} \right) = \mathbf{s}^{x}\left( \mathbf{n} \right) + \mathbf{s}^{x}\left( \mathbf{m} \right) \geq \mathbf{s}^{y}\left( \mathbf{n} \right) + \mathbf{s}^{y}\left( \mathbf{m} \right) = \mathbf{s}^{y}\left( \mathbf{n} + \mathbf{m} \right)\text{.}\]
Since this holds for all \(y \in \mathcal{X}\), we deduce that
\(x \in F\left( \mathbf{n} + \mathbf{m} \right)\). Since this works for
any
\(x \in F\left( \mathbf{n} \right) \cap F\left( \mathbf{m} \right)\), we
conclude that
\(F\left( \mathbf{n} \right) \cap F\left( \mathbf{m} \right) \subseteq F\left( \mathbf{n} + \mathbf{m} \right)\).

On the other hand, suppose
\(y \in \mathcal{X} \smallsetminus F\left( \mathbf{n} \right) \cap F\left( \mathbf{m} \right)\).
Then either \(y \notin F\left( \mathbf{n} \right)\) or
\(y \notin F\left( \mathbf{m} \right)\). Suppose without loss of
generality that \(y \notin F\left( \mathbf{n} \right)\). Let
\(x \in F\left( \mathbf{n} \right) \cap F\left( \mathbf{m} \right)\).
Then
\(\mathbf{s}^{y}\left( \mathbf{n} \right) < \mathbf{s}^{x}\left( \mathbf{n} \right)\),
while
\(\mathbf{s}^{y}\left( \mathbf{m} \right) \leq \mathbf{s}^{x}\left( \mathbf{m} \right)\).
Thus,
\[\mathbf{s}^{y}\left( \mathbf{n} + \mathbf{m} \right) = \mathbf{s}^{y}\left( \mathbf{n} \right) + \mathbf{s}^{y}\left( \mathbf{m} \right) < \mathbf{s}^{x}\left( \mathbf{n} \right) + \mathbf{s}^{x}\left( \mathbf{m} \right) = \mathbf{s}^{x}\left( \mathbf{n} + \mathbf{m} \right)\text{.}\]
Thus, \(y \notin F\left( \mathbf{n} + \mathbf{m} \right)\). Since this
works for any
\(y \in \mathcal{X} \smallsetminus F\left( \mathbf{n} \right) \cap F\left( \mathbf{m} \right)\),
we conclude that
\(F\left( \mathbf{n} \right) \cap F\left( \mathbf{m} \right) \supseteq F\left( \mathbf{n} + \mathbf{m} \right)\).

This proves ``\(\Longleftarrow\)''. Meanwhile, ``\(\Longrightarrow\)''
is an immediate consequence of Lemma~C.1, which we prove next.
□~Theorem~1

Let
\(\mathbb{Z}^{\langle\mathcal{V}\rangle} ≔ \left\{ \mathbf{n} \in \mathbb{Z}^{\mathcal{V}};\left\| \mathbf{n} \right\|_{} < \infty \right\}\);
then \(\mathbb{Z}^{\langle\mathcal{V}\rangle}\) is a torsion-free
abelian group. Let \(\mathcal{R}\) be another abelian group. There is a
bijective correspondence (actually, a group isomorphism) between
\(\mathcal{R}^{\mathcal{V}}\) and the set of group homomorphisms from
\(\mathbb{Z}^{\langle\mathcal{V}\rangle}\) into \(\mathcal{R}\). To see
this, note that any vector
\(\mathbf{b} = \left( b_{v} \right)_{v \in \mathcal{V}} \in \mathcal{R}^{\mathcal{V}}\)
defines a group homomorphism
\(\left. \mathbf{b}^{\ast}:\mathbb{Z}^{\langle\mathcal{V}\rangle}\longrightarrow\mathcal{R} \right.\)
by setting
\(\mathbf{b}^{\ast}\left( \mathbf{n} \right) ≔ \sum_{v \in \mathcal{V}}n_{v}b_{v}\)
for all \(\mathbf{n} \in \mathbb{Z}^{\langle\mathcal{V}\rangle}\).
Conversely, for all \(v \in \mathcal{V}\), let
\(\mathbf{1}_{v} \in \mathbb{Z}^{\langle\mathcal{V}\rangle}\) be the
vector with a 1 in the \(v\) coordinate, and 0 in all other coordinates.
Then given any homomorphism
\(\left. \beta:\mathbb{Z}^{\langle\mathcal{V}\rangle}\longrightarrow\mathcal{R} \right.\),
we can define a vector
\(\mathbf{b} = \left( b_{v} \right)_{v \in \mathcal{V}} \in \mathcal{R}^{\mathcal{V}}\)
by setting \(b_{v} ≔ \beta\left( \mathbf{1}_{v} \right)\) for all
\(v \in \mathcal{V}\). We then have \(\mathbf{b}^{\ast} = \beta\).

Thus, we can equivalently define a balance system
\(\left\{ \mathbf{b}^{x,y} \right\}_{x,y \in \mathcal{X}}\) or score
system \(\left\{ \mathbf{s}^{x} \right\}_{x \in \mathcal{X}}\) as a
collection of \(\mathcal{R}\)-valued group homomorphisms on
\(\mathbb{Z}^{\langle\mathcal{V}\rangle}\), rather than as a collection
of vectors in \(\mathcal{R}^{\mathcal{V}}\). We will make use of this
fact repeatedly in what follows. For simplicity, we will abuse notation
by using the same symbol (e.g.~`\(\mathbf{b}^{x,y}\)') to represent both
a vector in \(\mathcal{R}^{\mathcal{V}}\) and its associated group
homomorphism (i.e.~we will not write the `\(\ast\)'). This notational
convention should be kept in mind when reading the next proof.

Recall that a balance system \(\mathsf{B}\) is perfect on the domain
\(\mathcal{D}\) if, for any \(\mathbf{d} \in \mathcal{D}\), any
\(x \in F_{\mathsf{B}}\left( \mathbf{d} \right)\) and any
\(y \in \mathcal{X} \smallsetminus F_{\mathsf{B}}\left( \mathbf{d} \right)\),
we have \(\mathbf{b}^{x,y}\left( \mathbf{d} \right) > 0\). Not every
balance system is perfect (see Section~3). However, in addition to
completing the proof of Theorem~1, the next result implies that every
balance rule has a perfect balance system. Lemma~C.1

Let \(\mathcal{D} \subset \mathbb{N}^{\langle\mathcal{V}\rangle}\) be a
domain. If \(\left. F:\mathcal{D}\rightrightarrows\mathcal{X} \right.\)
satisfies reinforcement, then \(F = F_{\mathsf{B}}\) for some perfect
balance system \(\mathsf{B}\).

Proof

For all \(x \in \mathcal{X}\), let
\(\mathcal{C}_{x} ≔ \left\{ \mathbf{d} \in \mathcal{D};x \in F\left( \mathbf{d} \right) \right\}\).
Then \(\mathcal{C}_{x}\) is a preorder conoid in the abelian group
\(\mathbb{Z}^{\langle\mathcal{V}\rangle}\). (We have
\(\mathbf{0} \in \mathcal{C}_{x}\) because
\(\mathbf{0} \in \mathcal{D}\) and
\(F\left( \mathbf{0} \right) = \mathcal{X}\) by definition. Meanwhile,
\(\mathcal{C}_{x}\) is closed under addition because \(F\) satisfies
reinforcement.)

For any \(x,y \in \mathcal{X}\), let
\(\mathcal{P}_{x,y} ≔ \left\{ \mathbf{c}_{x} - \mathbf{c}_{y};\mathbf{c}_{x} \in \mathcal{C}_{x}\text{~and~}\mathbf{c}_{y} \in \mathcal{C}_{y} \right\}\).
Then \(\mathcal{P}_{x,y}\) is a preorder conoid. (\(\mathcal{P}_{x,y}\)
is closed under addition because \(\mathcal{C}_{x}\) and
\(\mathcal{C}_{y}\) are closed under addition.)~Note that
\(\mathcal{P}_{y,x} = - \mathcal{P}_{x,y}\).

Let \(( \succeq )\) be the homogeneous preorder on
\(\mathbb{Z}^{\langle\mathcal{V}\rangle}\) defined by
\(\mathcal{P}_{x,y}\) via formula (B.1). Let
\(\mathcal{O}_{x,y} ≔ \left\{ \mathbf{z} \in \mathbb{Z}^{\langle\mathcal{V}\rangle};\mathbf{z} \approx \mathbf{0} \right\}\),
and let
\(\mathcal{R}_{x,y} ≔ \mathbb{Z}^{\langle\mathcal{V}\rangle}/\mathcal{O}_{x,y}\),
with the homogeneous partial order \(( \succ )\) described in
Lemma~B.1(b). Use the Homogeneous Szpilrajn Lemma to extend
\(( \succ )\) to a homogeneous linear order
\(\left( \underset{x,y}{>} \right)\) on \(\mathcal{R}_{x,y}\). Let
\(\left. \mathbf{b}^{x,y}:\mathbb{Z}^{\langle\mathcal{V}\rangle}\longrightarrow\mathcal{R}_{x,y} \right.\)
be the quotient map
(i.e.~\(\mathbf{b}^{x,y}\left( \mathbf{z} \right) ≔ \mathbf{z} + \mathcal{O}_{x,y}\)
for all \(\mathbf{z} \in \mathbb{Z}^{\langle\mathcal{V}\rangle}\)). Note
that \(\mathcal{O}_{y,x} = \mathcal{O}_{x,y}\), so
\(\mathcal{R}_{y,x} = \mathcal{R}_{x,y}\) as groups, and
\(\mathbf{b}^{y,x} = \mathbf{b}^{x,y}\). However,
\(\left( \underset{x,y}{>} \right)\) is the negative ordering to
\(\left( \underset{y,x}{>} \right)\)
(i.e.~\(\left. \left( r\underset{x,y}{>}r' \right)\;\Longleftrightarrow\;\left( - r\underset{y,x}{>} - r' \right) \right.\)).
Thus, without loss of generality, we can redefine
\(\left( \underset{y,x}{>} \right)\) to be identical with
\(\left( \underset{x,y}{>} \right)\), and redefine \(\mathbf{b}^{y,x}\)
to be \(- \mathbf{b}^{x,y}\) (or vice versa; it does not matter whether
we reverse \(\mathbf{b}^{x,y}\) or \(\mathbf{b}^{y,x}\), as long as we
reverse only one of them). Finally, define \(\mathbf{b}^{x,x} ≔ 0\) for
all \(x \in \mathcal{X}\). Claim~C.1.1

Let \(\mathbf{d} \in \mathcal{D}\) and let
\(x \in F\left( \mathbf{d} \right)\) . Then:(a)

\(\mathbf{b}^{x,y}\left( \mathbf{d} \right) \geq 0\) for all
\(y \in \mathcal{X}\) (~hence,
\(x \in F_{\mathsf{B}}\left( \mathbf{d} \right)\) ).

(b)

Furthermore, if \(y \notin F\left( \mathbf{d} \right)\), then
\(\mathbf{b}^{x,y}\left( \mathbf{d} \right) > 0\) (~hence,
\(y \notin F_{\mathsf{B}}\left( \mathbf{d} \right)\) ).

Proof

(a)

If \(x \in F\left( \mathbf{d} \right)\), then
\(\mathbf{d} \in \mathcal{C}_{x}\). Meanwhile
\(\mathbf{0} \in \mathcal{C}_{y}\); thus,
\(\mathbf{d} = \mathbf{d} - \mathbf{0} \in \mathcal{P}_{x,y}\), so
\(\mathbf{d} \succeq \mathbf{0}\), so
\(\mathbf{b}^{x,y}\left( \mathbf{d} \right) \geq 0\).

(b)

(by contradiction) Suppose
\(\mathbf{b}^{x,y}\left( \mathbf{d} \right) = 0\). Then
\(\mathbf{d} \in \mathcal{O}_{x,y}\). Thus,
\(\mathbf{d} \preceq \mathbf{0}\), so
\(- \mathbf{d} \succeq \mathbf{0}\), so
\(- \mathbf{d} \in \mathcal{P}_{x,y}\). Thus,
\(- \mathbf{d} = \mathbf{c}_{x} - \mathbf{c}_{y}\) for some
\(\mathbf{c}_{x} \in \mathcal{C}_{x}\) and
\(\mathbf{c}_{y} \in \mathcal{C}_{y}\). But then
\(\mathbf{c}_{y} = \mathbf{c}_{x} + \mathbf{d}\). Now,
\(x \in F\left( \mathbf{d} \right)\) and
\(x \in F\left( \mathbf{c}_{x} \right)\), so
\(F\left( \mathbf{d} \right) \cap F\left( \mathbf{c}_{x} \right) \neq \varnothing\);
thus
\(F\left( \mathbf{c}_{y} \right) = F\left( \mathbf{d} + \mathbf{c}_{x} \right) = F\left( \mathbf{d} \right) \cap F\left( \mathbf{c}_{x} \right)\),
by reinforcement. But \(y \notin F\left( \mathbf{d} \right)\) by
hypothesis, so
\(y \notin F\left( \mathbf{d} \right) \cap F\left( \mathbf{c}_{x} \right)\),
so \(y \notin F\left( \mathbf{c}_{y} \right)\). But this contradicts the
fact that \(\mathbf{c}_{y} \in \mathcal{C}_{y}\). By contradiction, we
cannot have \(\mathbf{b}^{x,y}\left( \mathbf{d} \right) = 0\). Thus,
\(\mathbf{b}^{x,y}\left( \mathbf{d} \right) > 0\). ♢Claim~C.1.1

For all \(\mathbf{d} \in \mathcal{D}\), Claim~C.1.1(a) implies that
\(F\left( \mathbf{d} \right) \subseteq F_{\mathsf{B}}\left( \mathbf{d} \right)\),
while Claim~C.1.1(b) implies that
\(F\left( \mathbf{d} \right) \supseteq F_{\mathsf{B}}\left( \mathbf{d} \right)\).
We conclude that \(F = F_{\mathsf{B}}\). Furthermore, Claim~C.1.1(b)
implies that the balance system \(\mathsf{B}\) is perfect.

Now we must construct a single linearly ordered abelian group
\(\mathcal{R}\) and a collection of functions
\(\left. {\widetilde{\mathbf{b}}}_{x,y}:\mathcal{V}\longrightarrow\mathcal{R} \right.\)
(for all \(x,y \in \mathcal{X}\)) such that
\(F_{\mathsf{B}} = F_{\widetilde{\mathsf{B}}}\). Let
\(\mathcal{R} ≔ \prod_{x,y \in \mathcal{X}}\mathcal{R}_{x,y}\) with the
coordinatewise dominance order it gets from the linear orders on the
components. Then extend this to a homogeneous linear order on
\(\mathcal{R}\) using the Homogeneous Szpilrajn Lemma.□ Lemma~C.1

The proof of Theorem~2 depends upon Proposition~1, Proposition~2, so it
is deferred.

\protect\phantomsection\label{s000040}
Towards the proof of Proposition~1

Lemma~C.2

Let \(\mathbf{b} \in \mathcal{R}^{\mathcal{V}}\).(a)

For any \(\mathbf{n} \in \mathbb{N}^{\langle\mathcal{V}\rangle}\) and
\(\pi \in \Pi_{\mathcal{V}}\), we have
\(\left( \mathbf{b}\pi \right)\left( \mathbf{n} \right) = \mathbf{b}\left( \pi\left( \mathbf{n} \right) \right)\).

(b)

For any \(\pi,\phi \in \Pi_{\mathcal{V}}\), we have
\(\left( \mathbf{b}\,\pi \right)\phi = \mathbf{b}(\pi\phi)\).

Proof

(a)

\(\left( \mathbf{b}\pi \right)\left( \mathbf{n} \right) = \sum_{v \in \mathcal{V}}n_{v}\left( \mathbf{b}\pi \right)_{v} = \sum_{v \in \mathcal{V}}n_{v}b_{\pi{(v)}}\underset{( \ast )}{=}\sum_{v' \in \mathcal{V}}n_{\pi^{- 1}{(v')}}b_{v'} = \sum_{v \in \mathcal{V}}\pi\left( \mathbf{n} \right)_{v'}b_{v'} = \mathbf{b}\left( \pi\left( \mathbf{n} \right) \right)\).
Here, \(( \ast )\) is the change of variables \(v' ≔ \pi(v)\).

(b)

For any \(v \in \mathcal{V}\), if \(w ≔ \phi(v)\), then
\(\left( \left( \mathbf{b}\pi \right)\phi \right)_{v} = \left( \mathbf{b}\pi \right)_{\phi{(v)}} = \left( \mathbf{b}\pi \right)_{w} = b_{\pi{(w)}} = b_{\pi\phi{(v)}} = \left( \mathbf{b}\,(\pi\phi) \right)_{v}\),
as desired.□ Lemma~C.2

Proof of Proposition~1

``\(\Longleftarrow\)'' Let \(\mathsf{S}\) be a \(\nu\)-neutral score
system. Let \(\pi \in \Pi_{\mathcal{X}}'\), let
\(\widetilde{\pi} ≔ \nu(\pi) \in \Pi_{\mathcal{V}}\), and let
\(\mathbf{d} \in \mathcal{D}\). We must show that
\(F_{\mathsf{S}}\left\lbrack \widetilde{\pi}\left( \mathbf{d} \right) \right\rbrack = \pi\left\lbrack F_{\mathsf{S}}\left( \mathbf{d} \right) \right\rbrack\).
To see this, observe that, for any \(y \in \mathcal{X}\), if
\(x = \pi(y)\), then
(C.1)\[\mathbf{s}^{x}\left\lbrack \widetilde{\pi}\left( \mathbf{d} \right) \right\rbrack\underset{( \ast )}{=}\left( \mathbf{s}^{x}\widetilde{\pi} \right)\left( \mathbf{d} \right)\underset{( \dagger )}{=}\mathbf{s}^{y}\left( \mathbf{d} \right)\text{,}\]
where \(( \ast )\) is by Lemma~C.2(a) and \(( \dagger )\) because
\(\mathsf{B}\) is \(\nu\)-neutral. Thus,
\[\left. \left( y \in F_{\mathsf{B}}\left( \mathbf{d} \right) \right)\underset{( \dagger )}{\left. \Leftarrow = \Rightarrow \right.}\left( \mathbf{s}^{y}\left( \mathbf{d} \right) \geq \mathbf{s}^{z}\left( \mathbf{d} \right)\text{,~for~all~}z \in \mathcal{X} \right)\underset{( \ast )}{\left. \Leftarrow = \Rightarrow \right.}\left( \mathbf{s}^{x}\left\lbrack \widetilde{\pi}\left( \mathbf{d} \right) \right\rbrack \geq \mathbf{s}^{w}\left\lbrack \widetilde{\pi}\left( \mathbf{d} \right) \right\rbrack\text{,~for~all~}w \in \mathcal{X} \right)\underset{( \dagger )}{\left. \Leftarrow = \Rightarrow \right.}\left( x \in F_{\mathsf{B}}\left\lbrack \widetilde{\pi}\left( \mathbf{d} \right) \right\rbrack \right)\;\Longleftrightarrow\;\left( \pi(y) \in F_{\mathsf{B}}\left\lbrack \widetilde{\pi}\left( \mathbf{d} \right) \right\rbrack \right)\text{.} \right.\]
It follows that
\(F_{\mathsf{S}}\left\lbrack \widetilde{\pi}\left( \mathbf{d} \right) \right\rbrack = \pi\left\lbrack F_{\mathsf{S}}\left( \mathbf{d} \right) \right\rbrack\),
as desired. Here, \(( \dagger )\) is by definition of ``scoring rule'',
and \(( \ast )\) is by Eq. (C.1) (along with the change of variables
\(w = \pi(z)\)).

``\(\Longrightarrow\)'' For any \(\pi \in \Pi_{\mathcal{X}}'\), define
\(\widetilde{\pi} ≔ \nu(\pi)\) (an element of \(\Pi_{\mathcal{V}}\)).
Recall that \(\widetilde{\pi}\) acts on \(\mathcal{R}^{\mathcal{V}}\) by
coordinate permutations---i.e.~for any
\(\mathbf{r} = \left( r_{v} \right)_{v \in \mathcal{V}} \in \mathcal{R}^{\mathcal{V}}\),
we define
\(\mathbf{r}\widetilde{\pi} ≔ \left( r_{\widetilde{\pi}{(v)}} \right)_{v \in \mathcal{V}}\).

Now let
\(\mathsf{S} = \left\{ \mathbf{s}^{x} \right\}_{x \in \mathcal{X}}\) be
a score system such that \(F = F_{\mathsf{S}}\). Define the score system
\(\overline{\mathsf{S}} = \left\{ {\overline{\mathbf{s}}}^{x} \right\}_{x \in \mathcal{X}}\)
by setting
\[{\overline{\mathbf{s}}}^{x} ≔ \sum\limits_{\pi \in \Pi_{\mathcal{X}}'}\mathbf{s}^{\pi{(x)}}\,\widetilde{\pi}\text{,}\quad\text{for~all~}x \in \mathcal{X}\text{.}\]

Claim~1.1

\(F_{\overline{\mathsf{S}}} = F\).

Proof

`\(\supseteq\)' Let \(\mathbf{d} \in \mathcal{D}\) and
\(x \in \mathcal{X}\). Suppose \(x \in F\left( \mathbf{d} \right)\); we
will show that
\(x \in F_{\overline{\mathsf{S}}}\left( \mathbf{d} \right)\). For any
\(\pi \in \Pi_{\mathcal{X}}'\), neutrality implies that
\(\pi(x) \in F\left( \widetilde{\pi}\left( \mathbf{d} \right) \right)\).
Thus,
\(\mathbf{s}^{\pi{(x)}}\left( \widetilde{\pi}\left( \mathbf{d} \right) \right) \geq \mathbf{s}^{y}\left( \widetilde{\pi}\left( \mathbf{d} \right) \right)\)
for all \(y \in \mathcal{X}\) and \(\pi \in \Pi_{\mathcal{X}}'\).
Equivalently,
\(\mathbf{s}^{\pi{(x)}}\left( \widetilde{\pi}\left( \mathbf{d} \right) \right) \geq \mathbf{s}^{\pi{(z)}}\left( \widetilde{\pi}\left( \mathbf{d} \right) \right)\)
for all \(z \in \mathcal{X}\)and \(\pi \in \Pi_{\mathcal{X}}'\). Thus,
Lemma~C.2(a) says
\(\left( \mathbf{s}^{\pi{(x)}}\widetilde{\pi} \right)\left( \mathbf{d} \right) \geq \left( \mathbf{s}^{\pi{(z)}}\widetilde{\pi} \right)\left( \mathbf{d} \right)\)
for all \(z \in \mathcal{X}\) and \(\pi \in \Pi_{\mathcal{X}}'\). Thus,
\[{\overline{\mathbf{s}}}^{x}\left( \mathbf{d} \right) = \sum\limits_{\pi \in \Pi_{\mathcal{X}}'}\left( \mathbf{s}^{\pi{(x)}}\,\widetilde{\pi} \right)\left( \mathbf{d} \right) \geq \sum\limits_{\pi \in \underset{\mathcal{X}}{\overset{'}{\Pi}}}\left( \mathbf{s}^{\pi{(z)}}\,\widetilde{\pi} \right)\left( \mathbf{d} \right) = {\overline{\mathbf{s}}}^{z}\left( \mathbf{d} \right)\text{,}\]
for all \(z \in \mathcal{X}\). Thus,
\(x \in F_{\overline{\mathsf{S}}}\left( \mathbf{d} \right)\), as
desired.

`\(\subseteq\)' (by contrapositive)~Now suppose
\(x \notin F\left( \mathbf{d} \right)\); we will show that
\(x \notin F_{\overline{\mathsf{S}}}\left( \mathbf{d} \right)\). Let
\(y \in F\left( \mathbf{d} \right)\). Then for any
\(\pi \in \Pi_{\mathcal{X}}'\), neutrality implies that
\(\pi(x) \notin F\left( \widetilde{\pi}\left( \mathbf{d} \right) \right)\)
and
\(\pi(y) \in F\left( \widetilde{\pi}\left( \mathbf{d} \right) \right)\).
Thus,
\(\mathbf{s}^{\pi{(x)}}\left( \widetilde{\pi}\left( \mathbf{d} \right) \right) < \mathbf{s}^{\pi{(y)}}\left( \widetilde{\pi}\left( \mathbf{d} \right) \right)\)
for all \(\pi \in \Pi_{\mathcal{X}}'\). Thus, Lemma~C.2(a) says
\(\left( \mathbf{s}^{\pi{(x)}}\widetilde{\pi} \right)\left( \mathbf{d} \right) < \left( \mathbf{s}^{\pi{(y)}}\widetilde{\pi} \right)\left( \mathbf{d} \right)\)
for all \(\pi \in \Pi_{\mathcal{X}}'\). Thus,
\[{\overline{\mathbf{s}}}^{x}\left( \mathbf{d} \right) = \sum\limits_{\pi \in \Pi_{\mathcal{X}}'}\left( \mathbf{s}^{\pi{(x)}}\,\widetilde{\pi} \right)\left( \mathbf{d} \right) < \sum\limits_{\pi \in \underset{\mathcal{X}}{\overset{'}{\Pi}}}\left( \mathbf{s}^{\pi{(y)}}\,\widetilde{\pi} \right)\left( \mathbf{d} \right) = {\overline{\mathbf{s}}}^{y}\left( \mathbf{d} \right)\text{.}\]
Thus, \(x \notin F_{\overline{\mathsf{S}}}\left( \mathbf{d} \right)\),
as desired. ♢~Claim~1.1

It remains to check that the system
\(\overline{\mathsf{S}} ≔ \left\{ {\overline{\mathbf{s}}}^{x} \right\}_{x \in \mathcal{X}}\)
is \(\nu\)-neutral. To see this, let \(x,y \in \mathcal{X}\), let
\(\phi \in \Pi_{\mathcal{X}}'\), and suppose \(\phi(y) = x\). Then
\[{\overline{\mathbf{s}}}^{x}\,\widetilde{\phi} = \left( \sum\limits_{\pi \in \Pi_{\mathcal{X}}'}\mathbf{s}^{\pi{(x)}}\,\widetilde{\pi} \right)\widetilde{\phi}\underset{( \ast )}{=}\sum\limits_{\pi \in \Pi_{\mathcal{X}}'}\mathbf{s}^{\pi{(x)}}\,\left( \widetilde{\pi}\widetilde{\phi} \right)\underset{(\diamondsuit)}{=}\sum\limits_{\pi \in \underset{\mathcal{X}}{\overset{'}{\Pi}}}\mathbf{s}^{\pi{(x)}}\,\widetilde{\pi\phi}\underset{( \dagger )}{=}\sum\limits_{\pi' \in \underset{\mathcal{X}}{\overset{'}{\Pi}}}\mathbf{s}^{\pi'\phi^{- 1}{(x)}}\,{\widetilde{\pi}}' = \sum\limits_{\pi' \in \underset{\mathcal{X}}{\overset{'}{\Pi}}}\mathbf{s}^{\pi'{(y)}}\,{\widetilde{\pi}}' = {\overline{\mathbf{s}}}^{y}\text{,}\]
as desired. Here, \(( \ast )\) is by linearity and Lemma~C.2(b),
\((\diamondsuit)\) is because \(\nu\) is a group homomorphism, and
\(( \dagger )\) is the change of variables \(\pi' ≔ \pi\phi\).□
Proposition~1

\protect\phantomsection\label{s000045}
Towards the proof of Proposition~2

Lemma~C.3

Let \(\mathcal{X}\) be a finite set, let \(\mathcal{V}\) be an arbitrary
set, let
\(\mathcal{D} \subseteq \mathbb{N}^{\langle\mathcal{V}\rangle}\) be a
domain, and let
\(\left. \nu:\Pi_{\mathcal{X}}'\longrightarrow\Pi_{\mathcal{V}} \right.\)
be a homomorphism. Let \(\mathcal{R}\) be a linearly ordered abelian
group, and let \(\widetilde{\mathsf{B}}\) be a \(\mathcal{R}\)-valued
perfect balance system on \(\left( \mathcal{X},\mathcal{V} \right)\) .
The balance rule
\(\left. F_{\widetilde{\mathsf{B}}}:\mathcal{D}\rightrightarrows\mathcal{X} \right.\)
is \(\nu\)-neutral if and only if
\(F_{\widetilde{\mathsf{B}}} = F_{\mathsf{B}}\) for some
\(\mathcal{R}\)-valued, \(\nu\)-neutral, perfect balance system
\(\mathsf{B}\) on \(\left( \mathcal{X},\mathcal{V} \right)\).

Proof

``\(\Longleftarrow\)'' Let \(\mathsf{B}\) be a \(\nu\)-neutral balance
system. Let \(\mathbf{d} \in \mathcal{D}\) and let
\(\pi \in \Pi_{\mathcal{X}}'\). We must show that
\(F_{\mathsf{B}}\left( \widetilde{\pi}\left( \mathbf{d} \right) \right) = \pi\left( F_{\mathsf{B}}\left( \mathbf{d} \right) \right)\).

``\(\supseteq\)'' Let \(x \in F_{\mathsf{B}}\left( \mathbf{d} \right)\),
and let \(y = \pi(x)\). For all \(z \in \mathcal{X}\), if
\(w = \pi^{- 1}(z)\), then
(C.2)\[\mathbf{b}^{y,z}\left( \widetilde{\pi}\left( \mathbf{d} \right) \right)\underset{( \ast )}{=}\left( \mathbf{b}^{y,z}\widetilde{\pi} \right)\left( \mathbf{d} \right)\underset{( \dagger )}{=}\mathbf{b}^{x,w}\left( \mathbf{d} \right)\underset{(\diamondsuit)}{\geq}0\text{.}\]
Here, \(( \ast )\) is by Lemma~C.2(a), \(( \dagger )\) is because
\(\mathsf{B}\) is \(\nu\)-neutral, and \((\diamondsuit)\) is because
\(x \in F_{\mathsf{B}}\left( \mathbf{d} \right)\). Since this holds for
all \(z \in \mathcal{X}\), we conclude that
\(y \in F_{\mathsf{B}}\left( \widetilde{\pi}\left( \mathbf{d} \right) \right)\).

``\(\subseteq\)'' (by contrapositive) Now suppose
\(x \notin F_{\mathsf{B}}\left( \mathbf{d} \right)\). Then there is some
\(w \in \mathcal{X}\) such that
\(\mathbf{b}^{x,w}\left( \mathbf{d} \right) < 0\). Thus, if
\(y = \pi(x)\) and \(z = \pi(w)\), then the inequality
`\(\underset{(\diamondsuit)}{\geq}\)' in Eq. (C.2) changes to `\(<\)';
thus,
\(y \notin F_{\mathsf{B}}\left( \widetilde{\pi}\left( \mathbf{d} \right) \right)\).

Thus
\(F_{\mathsf{B}}\left( \widetilde{\pi}\left( \mathbf{d} \right) \right) = \pi\left( F_{\mathsf{B}}\left( \mathbf{d} \right) \right)\).
This holds for all \(\mathbf{d} \in \mathcal{D}\) and
\(\pi \in \Pi_{\mathcal{X}}'\). Thus, \(F_{\mathsf{B}}\) is
\(\nu\)-neutral.

``\(\Longrightarrow\)'' Fix \(x,y \in \mathcal{X}\) with \(x \neq y\).
Double transitivity yields some \(\alpha \in \Pi_{\mathcal{X}}'\) such
that \(\alpha(x) = y\) and \(\alpha(y) = x\). Let
\(\widetilde{\alpha} ≔ \nu(\alpha) \in \Pi_{\mathcal{V}}\). Let \(N\) be
the order of \(\alpha\) (so that \(\alpha^{N} = {Id}_{\mathcal{X}}\) and
\({\widetilde{\alpha}}^{N} = {Id}_{\mathcal{V}}\)). Then \(N\) is
finite, because \(\mathcal{X}\) is finite. Also, \(N\) is even, because
\(\alpha(x) = y\) and \(\alpha^{2}(x) = x\). Let \(M = N/2\). Define
(C.3)\[{\overline{\mathbf{b}}}^{x,y} ≔ \sum\limits_{m = 0}^{M - 1}\left( {\widetilde{\mathbf{b}}}^{x,y}{\widetilde{\alpha}}^{2m} - {\widetilde{\mathbf{b}}}^{x,y}{\widetilde{\alpha}}^{2m + 1} \right)\text{.}\]

Claim~C.3.1

\({\overline{\mathbf{b}}}^{x,y}\widetilde{\alpha} = - {\overline{\mathbf{b}}}^{x,y}\).

Proof

We have
\[{\overline{\mathbf{b}}}^{x,y}\widetilde{\alpha}\underset{}{=}\left( \sum\limits_{m = 0}^{M - 1}{\widetilde{\mathbf{b}}}^{x,y}{\widetilde{\alpha}}^{2m} - \sum\limits_{m = 0}^{M - 1}{\widetilde{\mathbf{b}}}^{x,y}{\widetilde{\alpha}}^{2m + 1} \right)\widetilde{\alpha} = \sum\limits_{m = 0}^{M - 1}{\widetilde{\mathbf{b}}}^{x,y}{\widetilde{\alpha}}^{2m + 1} - \sum\limits_{m = 0}^{M - 1}{\widetilde{\mathbf{b}}}^{x,y}{\widetilde{\alpha}}^{2m + 2}\underset{( \ast )}{=}\sum\limits_{m = 0}^{M - 1}{\widetilde{\mathbf{b}}}^{x,y}{\widetilde{\alpha}}^{2m + 1} - \sum\limits_{m' = 1}^{M - 1}{\widetilde{\mathbf{b}}}^{x,y}{\widetilde{\alpha}}^{2m'} - {\widetilde{\mathbf{b}}}^{x,y}{\widetilde{\alpha}}^{2M}\underset{( \dagger )}{=}\sum\limits_{m = 0}^{M - 1}{\widetilde{\mathbf{b}}}^{x,y}{\widetilde{\alpha}}^{2m + 1} - \sum\limits_{m' = 1}^{M - 1}{\widetilde{\mathbf{b}}}^{x,y}{\widetilde{\alpha}}^{2m'} - {\widetilde{\mathbf{b}}}^{x,y} = \sum\limits_{m = 0}^{M - 1}{\widetilde{\mathbf{b}}}^{x,y}{\widetilde{\alpha}}^{2m + 1} - \sum\limits_{m' = 0}^{M - 1}{\widetilde{\mathbf{b}}}^{x,y}{\widetilde{\alpha}}^{2m'}\underset{}{=} - {\overline{\mathbf{b}}}^{x,y}\text{,}\]
as desired. Here, \(( \ast )\) is the change of variables
\(m' ≔ m + 1\), while \(( \dagger )\) is because
\({\widetilde{\alpha}}^{2M} = {\widetilde{\alpha}}^{N} = {Id}_{\mathcal{V}}\).
♢Claim~C.3.1

Recall that
\(\left. F_{\widetilde{\mathsf{B}}}:\mathcal{D}\rightrightarrows\mathcal{X} \right.\)
is the balance rule defined by the perfect balance system
\(\widetilde{\mathsf{B}}\).

Claim~C.3.2

Let \(\mathbf{d} \in \mathcal{D}\).(a)

If \(x \in F_{\widetilde{\mathsf{B}}}\left( \mathbf{d} \right)\), then
\({\overline{\mathbf{b}}}^{x,y}\left( \mathbf{d} \right) \geq 0\).

(b)

If \(x \notin F_{\widetilde{\mathsf{B}}}\left( \mathbf{d} \right)\), but
\(y \in F_{\widetilde{\mathsf{B}}}\left( \mathbf{d} \right)\), then
\({\overline{\mathbf{b}}}^{x,y}\left( \mathbf{d} \right) < 0\).

Proof

(a)

If \(x \in F_{\widetilde{\mathsf{B}}}\left( \mathbf{d} \right)\), then
for all \(m \in \lbrack 0\ldots M)\), the neutrality of
\(F_{\widetilde{\mathsf{B}}}\) implies that
\(x \in F_{\widetilde{\mathsf{B}}}\left( {\widetilde{\alpha}}^{2m}\left( \mathbf{d} \right) \right)\),
while
\(y \in F_{\widetilde{\mathsf{B}}}\left( {\widetilde{\alpha}}^{2m + 1}\left( \mathbf{d} \right) \right)\)
(because \(\alpha^{2m}(x) = x\), while \(\alpha^{2m + 1}(x) = y\)).
Thus,
\({\widetilde{\mathbf{b}}}^{x,y}\left( {\widetilde{\alpha}}^{2m}\left( \mathbf{d} \right) \right) \geq 0\),
while
\({\widetilde{\mathbf{b}}}^{x,y}\left( {\widetilde{\alpha}}^{2m + 1}\left( \mathbf{d} \right) \right) \leq 0\).
Thus, Lemma~C.2(a) says
\(\left( {\widetilde{\mathbf{b}}}^{x,y}{\widetilde{\alpha}}^{2m} \right)\left( \mathbf{d} \right) \geq 0\),
while
\(\left( {\widetilde{\mathbf{b}}}^{x,y}{\widetilde{\alpha}}^{2m + 1} \right)\left( \mathbf{d} \right) \leq 0\).
Thus,\[{\overline{\mathbf{b}}}^{x,y}\left( \mathbf{d} \right)\underset{}{=}\sum\limits_{m = 0}^{M - 1}\left( \left( {\widetilde{\mathbf{b}}}^{x,y}{\widetilde{\alpha}}^{2m} \right)\left( \mathbf{d} \right) - \left( {\widetilde{\mathbf{b}}}^{x,y}{\widetilde{\alpha}}^{2m + 1} \right)\left( \mathbf{d} \right) \right) \geq 0\text{.}\]

(b)

If \(x \notin F_{\widetilde{\mathsf{B}}}\left( \mathbf{d} \right)\) and
\(y \in F_{\widetilde{\mathsf{B}}}\left( \mathbf{d} \right)\), then for
all \(m \in \lbrack 0\ldots M)\), neutrality implies that
\(x \notin F_{\widetilde{\mathsf{B}}}\left( {\widetilde{\alpha}}^{2m}\left( \mathbf{d} \right) \right)\)
and
\(y \in F_{\widetilde{\mathsf{B}}}\left( {\widetilde{\alpha}}^{2m}\left( \mathbf{d} \right) \right)\),
while
\(y \notin F_{\widetilde{\mathsf{B}}}\left( {\widetilde{\alpha}}^{2m + 1}\left( \mathbf{d} \right) \right)\)
and
\(x \in F_{\widetilde{\mathsf{B}}}\left( {\widetilde{\alpha}}^{2m + 1}\left( \mathbf{d} \right) \right)\)
(because \(\alpha^{2m}(x) = x\) and \(\alpha^{2m}(y) = y\), while
\(\alpha^{2m + 1}(x) = y\) and \(\alpha^{2m + 1}(y) = x\)). Thus,
\({\widetilde{\mathbf{b}}}^{x,y}\left( {\widetilde{\alpha}}^{2m}\left( \mathbf{d} \right) \right) < 0\),
while
\({\widetilde{\mathbf{b}}}^{x,y}\left( {\widetilde{\alpha}}^{2m + 1}\left( \mathbf{d} \right) \right) > 0\),
because \(\widetilde{\mathsf{B}}\) is perfect. Thus, Lemma~C.2(a) says
\(\left( {\widetilde{\mathbf{b}}}^{x,y}{\widetilde{\alpha}}^{2m} \right)\left( \mathbf{d} \right) < 0\),
while
\(\left( {\widetilde{\mathbf{b}}}^{x,y}{\widetilde{\alpha}}^{2m + 1} \right)\left( \mathbf{d} \right) > 0\).
Thus,
\[{\overline{\mathbf{b}}}^{x,y}\left( \mathbf{d} \right)\underset{}{=}\sum\limits_{m = 0}^{M - 1}\left( \left( {\widetilde{\mathbf{b}}}^{x,y}{\widetilde{\alpha}}^{2m} \right)\left( \mathbf{d} \right) - \left( {\widetilde{\mathbf{b}}}^{x,y}{\widetilde{\alpha}}^{2m + 1} \right)\left( \mathbf{d} \right) \right) < 0\text{,}\]
because every summand is negative. ♢Claim~C.3.2

Let
\(\Pi_{x,y} ≔ \left\{ \pi \in \Pi_{\mathcal{X}}';\pi(x) = x\text{~and~}\pi(y) = y \right\}\)
(a subgroup of \(\Pi_{\mathcal{X}}'\)). Then \(\Pi_{x,y}\) is finite
because \(\Pi_{\mathcal{X}}'\) is finite (and \(\Pi_{\mathcal{X}}'\) is
finite because \(\mathcal{X}\) is finite). Define
(C.4)\[\mathbf{b}^{x,y} ≔ \sum\limits_{\pi \in \Pi_{x,y}}\left( {\overline{\mathbf{b}}}^{x,y}\widetilde{\pi} \right)\text{.}\]

Claim~C.3.3

\(\mathbf{b}^{x,y}\widetilde{\alpha} = - \mathbf{b}^{x,y}\).

Proof

\[\mathbf{b}^{x,y}\widetilde{\alpha}\underset{}{=}\left( \sum\limits_{\pi \in \Pi_{x,y}}{\overline{\mathbf{b}}}^{x,y}\,\widetilde{\pi} \right)\,\widetilde{\alpha} = \sum\limits_{\pi \in \Pi_{x,y}}\left( {\overline{\mathbf{b}}}^{x,y}\,\widetilde{\pi} \right)\,\widetilde{\alpha}\underset{(\diamondsuit)}{=}\sum\limits_{\pi \in \Pi_{x,y}}{\overline{\mathbf{b}}}^{x,y}\,\widetilde{\pi\alpha} = \sum\limits_{\pi \in \Pi_{x,y}}{\overline{\mathbf{b}}}^{x,y}\widetilde{\left( \alpha\alpha^{-1}\pi\alpha \right)}\underset{(\diamondsuit)}{=}\sum\limits_{\pi \in \Pi_{x,y}}\left( {\overline{\mathbf{b}}}^{x,y}\,\widetilde{\alpha} \right)\widetilde{\alpha^{-1}\pi\alpha}\underset{( \ast )}{=}\sum\limits_{\pi' \in \Pi_{x,y}}\left( {\overline{\mathbf{b}}}^{x,y}\,\widetilde{\alpha} \right)\,{\widetilde{\pi}}'\underset{( \dagger )}{=}\sum\limits_{\pi' \in \Pi_{x,y}} - {\overline{\mathbf{b}}}^{x,y}\,{\widetilde{\pi}}' = - \sum\limits_{\pi' \in \Pi_{x,y}}{\overline{\mathbf{b}}}^{x,y}\,{\widetilde{\pi}}'\underset{}{=} - \mathbf{b}^{x,y}\text{.}\]
Here, both \((\diamondsuit)\) are by Lemma~C.2(b) (and the fact that
\(\nu\) is a homomorphism). Meanwhile, \(( \ast )\) is by the change of
variables \(\pi' ≔ \alpha\pi\alpha^{- 1}\) for all \(\pi \in \Pi_{x,y}\)
(observe that the map \(\left( \pi\mapsto\alpha\pi\alpha^{- 1} \right)\)
is a permutation of \(\Pi_{x,y}\), because \(\alpha(x) = y\) and
\(\alpha(y) = x\)). Finally, \(( \dagger )\) is by Claim~C.3.1.
♢~Claim~C.3.3

Claim~C.3.4

For any \(\phi \in \Pi_{x,y}\), we have
\(\mathbf{b}^{x,y}\,\widetilde{\phi} = \mathbf{b}^{x,y}\).

Proof

\(\mathbf{b}^{x,y}\,\widetilde{\phi}\underset{}{=}\left( \sum_{\pi \in \Pi_{x,y}}{\overline{\mathbf{b}}}^{x,y}\,\widetilde{\pi} \right)\widetilde{\phi}\underset{(\diamondsuit)}{=}\sum_{\pi \in \Pi_{x,y}}{\overline{\mathbf{b}}}^{x,y}\,\widetilde{\pi\phi}\underset{( \ast )}{=}\sum_{\pi' \in \Pi_{x,y}}{\overline{\mathbf{b}}}^{x,y}\,{\widetilde{\pi}}'\underset{}{=}\mathbf{b}^{x,y}\).

Here, \((\diamondsuit)\) is by linearity and Lemma~C.2(b), while
\(( \ast )\) is the change of variables \(\pi' ≔ \pi\phi\) (the map
\(\left( \pi\mapsto\pi\phi \right)\) is a permutation of \(\Pi_{x,y}\),
because \(\phi \in \Pi_{x,y}\)). ♢Claim~C.3.4

Now, for all \(w,z \in \mathcal{X}\), define
\(\Pi_{w,z} ≔ \left\{ \pi \in \Pi_{\mathcal{X}}';\pi(w) = x \right.\)
and \(\left. \pi(z) = y \right\}\); then
\(\Pi_{w,z} \neq \varnothing\) by the double-transitivity of
\(\Pi_{\mathcal{X}}'\). Define
\(\mathbf{b}^{w,z} ≔ \mathbf{b}^{x,y}\,{\widetilde{\pi}}_{w,z}\), for
any \(\pi_{w,z} \in \Pi_{w,z}\). Claim~C.3.5

\(\mathbf{b}^{w,z}\) is well-defined independent of the choice of
\(\pi_{w,z} \in \Pi_{w,z}\).

Proof

Suppose \(\pi_{w,z} \in \Pi_{w,z}\) and \(\phi_{w,z} \in \Pi_{w,z}\).
Then \(\pi_{w,z}\,\phi_{w,z}^{- 1} \in \Pi_{x,y}\). Thus,
(C.5)\[\left( {\overline{\mathbf{b}}}^{x,y}\,{\widetilde{\pi}}_{w,z} \right)\,{\widetilde{\phi}}_{w,z}^{- 1}\underset{( \ast )}{=}{\overline{\mathbf{b}}}^{x,y}\left( \widetilde{\pi_{w,z}\,\phi_{w,z}^{-1}} \right)\underset{( \dagger )}{=}{\overline{\mathbf{b}}}^{x,y}\text{,}\]
where \(( \ast )\) is by Lemma~C.2(b), and \(( \dagger )\) is by
Claim~C.3.4. Multiplying both ends of (C.5) by
\({\widetilde{\phi}}_{w,z}\), we obtain:
\({\overline{\mathbf{b}}}^{x,y}\,{\widetilde{\pi}}_{w,z} = {\overline{\mathbf{b}}}^{x,y}\,{\widetilde{\phi}}_{w,z}\),
as desired. ♢Claim~C.3.5

Claim~C.3.6

For any \(w,z,w',z' \in \mathcal{X}\) and
\(\phi \in \Pi_{\mathcal{X}}'\), if \(\phi^{- 1}(w) = w'\) and
\(\phi^{- 1}(z) = z'\), then
\(\mathbf{b}^{w,z}\widetilde{\phi} = \mathbf{b}^{w',z'}\).

Proof

Let \(\pi_{w,z} \in \Pi_{w,z}\). Then
\(\mathbf{b}^{w,z} = \mathbf{b}^{x,y}{\widetilde{\pi}}_{w,z}\) by
definition (and Claim~C.3.5). But \(\pi_{w,z}\phi \in \Pi_{w',z'}\)
(because \(\pi_{w,z}\phi\left( w' \right) = \pi_{w,z}(w) = x\) and
\(\pi_{w,z}\phi\left( z' \right) = \pi_{w,z}(z) = y\)). Thus,
\[\mathbf{b}^{w,z}\widetilde{\phi} = \left( \mathbf{b}^{x,y}\,{\widetilde{\pi}}_{w,z} \right)\,\widetilde{\phi}\underset{( \ast )}{=}\mathbf{b}^{x,y}\,\left( \widetilde{\pi_{w,z}\phi} \right)\underset{( \dagger )}{=}\mathbf{b}^{w',z'}\text{,}\]
as desired. Here, \(( \ast )\) is by Lemma~C.2(b) and \(( \dagger )\) is
by definition (and Claim~C.3.5). ♢~Claim~C.3.6

Claim~C.3.7

\(\mathbf{b}^{z,w} = - \mathbf{b}^{w,z}\) for all
\(w,z \in \mathcal{X}\).

Proof

Let \(\pi_{w,z} \in \Pi_{w,z}\). Then \(\alpha\pi_{w,z} \in \Pi_{z,w}\)
(because \(\alpha\pi_{w,z}(z) = \alpha(y) = x\), etc.). Thus,
\[\mathbf{b}^{z,w}\underset{( \ast )}{=}\mathbf{b}^{x,y}\,\left( \widetilde{\alpha\pi_{w,z}} \right)\underset{( \ddagger )}{=}\left( \mathbf{b}^{x,y}\widetilde{\alpha} \right){\widetilde{\pi}}_{w,z}\underset{(\diamondsuit)}{=} - \mathbf{b}^{x,y}{\widetilde{\pi}}_{w,z}\underset{( \dagger )}{=} - \mathbf{b}^{w,z}\text{.}\]
Here, \(( \ast )\) is by definition of \(\mathbf{b}^{z,w}\) (and
Claim~C.3.5). Next, \(( \ddagger )\) is by Lemma~C.2(b),
\((\diamondsuit)\) is by Claim~C.3.3, and \(( \dagger )\) is by
definition of \(\mathbf{b}^{w,z}\) (and Claim~C.3.5). ♢~Claim~C.3.7

Now define
\(\mathsf{B} ≔ \left\{ \mathbf{b}^{x,y} \right\}_{x,y \in \mathcal{X}}\).
Then \(\mathsf{B}\) is a balance system by Claim~C.3.7, and
\(\mathsf{B}\) is \(\nu\)-neutral by Claim~C.3.6. The
`\(\Longleftarrow\)' direction of the theorem (already established) then
implies that the balance rule \(F_{\mathsf{B}}\) is \(\nu\)-neutral.

Claim~C.3.8

\(F_{\mathsf{B}}\left( \mathbf{d} \right) = F_{\widetilde{\mathsf{B}}}\left( \mathbf{d} \right)\)
for all \(\mathbf{d} \in \mathcal{D}\) . Also, \(\mathsf{B}\) is a
perfect balance rule on \(\mathcal{D}\).

Proof

``\(\supseteq\)'' Let \(\mathbf{d} \in \mathcal{D}\) and let
\(w \in F_{\widetilde{\mathsf{B}}}\left( \mathbf{d} \right)\). Fix
\(z \in \mathcal{X}\) and let \(\pi_{w,z} \in \Pi_{w,z}\). Then
\(\pi(w) = x\) and \(\pi(z) = y\). Since \(F_{\widetilde{\mathsf{B}}}\)
is \(\nu\)-neutral, we have
\(F_{\widetilde{\mathsf{B}}}\left( {\widetilde{\pi}}_{w,z}\left( \mathbf{d} \right) \right) = \pi_{w,z}\left( F_{\widetilde{\mathsf{B}}}\left( \mathbf{d} \right) \right)\);
thus
\(x \in F_{\widetilde{\mathsf{B}}}\left( {\widetilde{\pi}}_{w,z}\left( \mathbf{d} \right) \right)\).
Now, for any \(\phi \in \Pi_{x,y}\), neutrality implies that
\(x \in F_{\widetilde{\mathsf{B}}}\left( \widetilde{\phi}\,{\widetilde{\pi}}_{w,z}\left( \mathbf{d} \right) \right)\);
thus, Claim~C.3.2(a) says that
\({\overline{\mathbf{b}}}^{x,y}\left( \widetilde{\phi}\,{\widetilde{\pi}}_{w,z}\left( \mathbf{d} \right) \right) \geq 0\).
But Lemma~C.2(a) says
\({\overline{\mathbf{b}}}^{x,y}\left( \widetilde{\phi}\,{\widetilde{\pi}}_{w,z}\left( \mathbf{d} \right) \right) = \left( {\overline{\mathbf{b}}}^{x,y}\widetilde{\phi} \right)\left( {\widetilde{\pi}}_{w,z}\left( \mathbf{d} \right) \right)\).
Thus, we have
(C.6)\[\left( {\overline{\mathbf{b}}}^{x,y}\,\widetilde{\phi} \right)\left( {\widetilde{\pi}}_{w,z}\left( \mathbf{d} \right) \right) \geq 0\text{,}\quad\text{for~all~}\phi \in \Pi_{x,y}\text{.}\](C.7)\[\text{Thus},\mathbf{b}^{w,z}\left( \mathbf{d} \right)\underset{( \ast )}{=}\left( \mathbf{b}^{x,y}{\widetilde{\pi}}_{w,z} \right)\left( \mathbf{d} \right)\underset{( \dagger )}{=}\mathbf{b}^{x,y}\left( {\widetilde{\pi}}_{w,z}\left( \mathbf{d} \right) \right)\underset{}{=}\sum\limits_{\phi \in \Pi_{x,y}}\left( {\overline{\mathbf{b}}}^{x,y}\,\widetilde{\phi} \right)\left( {\widetilde{\pi}}_{w,z}\left( \mathbf{d} \right) \right)\underset{}{\geq}0\text{.}\]
Here, \(( \ast )\) is by definition of \(\mathbf{b}^{w,z}\) (and
Claim~C.3.5), while \(( \dagger )\) is by Lemma~C.2(a).

Inequality (C.7) holds for any \(z \in \mathcal{X}\). Thus,
\(w \in F_{\mathsf{B}}\left( \mathbf{d} \right)\). ``\(\subseteq\)''

(by contrapositive)~Suppose
\(w \notin F_{\widetilde{\mathsf{B}}}\left( \mathbf{d} \right)\). Let
\(z \in F_{\widetilde{\mathsf{B}}}\left( \mathbf{d} \right)\). We will
show that \(\mathbf{b}^{w,z}\left( \mathbf{d} \right) < 0\), thereby
showing both that \(w \notin F_{\mathsf{B}}\left( \mathbf{d} \right)\)
and that \(\mathsf{B}\) is perfect.

If \(\pi_{w,z} \in \Pi_{w,z}\), then the neutrality of
\(F_{\widetilde{\mathsf{B}}}\) implies that
\(x \notin F_{\widetilde{\mathsf{B}}}\left( {\widetilde{\pi}}_{w,z}\left( \mathbf{d} \right) \right)\),
while
\(y \in F_{\widetilde{\mathsf{B}}}\left( {\widetilde{\pi}}_{w,z}\left( \mathbf{d} \right) \right)\).
Thus, for all \(\phi \in \Pi_{x,y}\), we have
\(x \notin F_{\widetilde{\mathsf{B}}}\left( \widetilde{\phi}\,{\widetilde{\pi}}_{w,z}\left( \mathbf{d} \right) \right)\),
while
\(y \in F_{\widetilde{\mathsf{B}}}\left( \widetilde{\phi}\,{\widetilde{\pi}}_{w,z}\left( \mathbf{d} \right) \right)\).
Thus, Claim~C.3.2(b) says that
\({\overline{\mathbf{b}}}^{x,y}\left( \widetilde{\phi}\,{\widetilde{\pi}}_{w,z}\left( \mathbf{d} \right) \right) < 0\);
hence Lemma~C.2(a) says that
\(\left( {\overline{\mathbf{b}}}^{x,y}\,\widetilde{\phi} \right)\left( {\widetilde{\pi}}_{w,z}\left( \mathbf{d} \right) \right) < 0\).
This holds for all \(\phi \in \Pi_{x,y}\); thus, an argument similar to
Eq. (C.7) yields \(\mathbf{b}^{w,z}\left( \mathbf{d} \right) < 0\).
Thus, \(w \notin F_{\mathsf{B}}\left( \mathbf{d} \right)\).
♢~Claim~C.3.8

The ``\(\Longrightarrow\)'' direction of the lemma now follows from
Claim~C.3.6, Claim~C.3.8. □ Lemma~C.3

Proof of Proposition~2

``\(\Longleftarrow\)'' follows immediately from
Lemma~C.3\(\left( \Leftarrow \right)\).

``\(\Longrightarrow\)''~If \(F\) is a balance rule and \(\mathcal{D}\)
is weakly additive, then \(F\) satisfies reinforcement (by Theorem~1).
Then \(F = F_{\widetilde{\mathsf{B}}}\) for some perfect balance system
\(\widetilde{\mathsf{B}}\) (by Lemma~C.1). Now apply
Lemma~C.3\(\left( \Rightarrow \right)\). □~Proposition~2

\protect\phantomsection\label{s000050}
Towards the proof of Theorem~2

Lemma~C.4

Let \(\mathcal{D} \subseteq \mathbb{N}^{\langle\mathcal{V}\rangle}\) be
a domain. A voting rule
\(\left. F:\mathcal{D}\rightrightarrows\mathcal{X} \right.\) is a
scoring rule if and only if \(F\) is a balance rule with a balance
system
\(\mathsf{B} = \left\{ \mathbf{b}^{x,y} \right\}_{x,y \in \mathcal{X}}\)
satisfying:(C.8)\[\mathbf{b}^{x,y}\left( \mathbf{d} \right) + \mathbf{b}^{y,z}\left( \mathbf{d} \right) = \mathbf{b}^{x,z}\left( \mathbf{d} \right)\text{,}\quad\ x,y,z \in \mathcal{X}\text{~~}\mathbf{d} \in \mathcal{D}\]

Proof

`\(\Longrightarrow\)' Let
\(\mathsf{S} = \left\{ \mathbf{s}^{x} \right\}_{x \in \mathcal{X}}\) be
an \(\mathcal{R}\)-valued score system on
\(\left( \mathcal{X},\mathcal{V} \right)\). For all
\(x,y \in \mathcal{X}\), define
\(\mathbf{b}^{x,y} ≔ \mathbf{s}^{x} - \mathbf{s}^{y} \in \mathcal{R}^{\mathcal{V}}\),
to obtain a balance system
\(\mathsf{B} ≔ \left\{ \mathbf{b}^{x,y} \right\}_{x,y \in \mathcal{X}}\).
Then
\(F_{\mathsf{B}}\left( \mathbf{n} \right) = F_{\mathsf{S}}\left( \mathbf{n} \right)\)
for all \(\mathbf{n} \in \mathbb{N}^{\langle\mathcal{V}\rangle}\) (as
observed in Example~1).

`\(\Longleftarrow\)' Let
\(\mathsf{B} = \left\{ \mathbf{b}^{x,y} \right\}_{x,y \in \mathcal{X}}\)
be a balance system on \(\mathcal{V}\) satisfying (C.8), and such that
\(F_{\mathsf{B}}\left( \mathbf{d} \right) = F\left( \mathbf{d} \right)\)
for all \(\mathbf{d} \in \mathcal{D}\). Fix \(o \in \mathcal{X}\), and
define \(\mathbf{s}^{o} ≔ \mathbf{0}\). Then for all
\(x \in \mathcal{X}\), define \(\mathbf{s}^{x} ≔ \mathbf{b}^{x,o}\).
This yields a score system
\(\mathsf{S} = \left\{ \mathbf{s}^{x} \right\}_{x \in \mathcal{X}}\) on
\(\mathcal{V}\). Let
\({\widetilde{\mathbf{b}}}^{x,y} ≔ \mathbf{s}^{x} - \mathbf{s}^{y}\) for
all \(x,y \in \mathcal{X}\). Thus, Example~1 implies that
\(F_{\mathsf{S}}\left( \mathbf{n} \right) = F_{\widetilde{\mathsf{B}}}\left( \mathbf{n} \right)\)
for all \(\mathbf{n} \in \mathbb{N}^{\langle\mathcal{V}\rangle}\). For
any \(\mathbf{d} \in \mathcal{D}\), and all \(x,y \in \mathcal{X}\), we
have
\[{\widetilde{\mathbf{b}}}^{x,y}\left( \mathbf{d} \right) = \mathbf{s}^{x}\left( \mathbf{d} \right) - \mathbf{s}^{y}\left( \mathbf{d} \right) = \mathbf{b}^{x,o}\left( \mathbf{d} \right) - \mathbf{b}^{y,o}\left( \mathbf{d} \right)\underset{( \ast )}{=}\mathbf{b}^{x,y}\left( \mathbf{d} \right)\text{,}\]
where \(( \ast )\) is by condition (C.8). Thus,
\(F_{\widetilde{\mathsf{B}}}\left( \mathbf{d} \right) = F_{\mathsf{B}}\left( \mathbf{d} \right)\)
for all \(\mathbf{d} \in \mathcal{D}\). But
\(F_{\mathsf{B}}\left( \mathbf{d} \right) = F\left( \mathbf{d} \right)\)
for all \(\mathbf{d} \in \mathcal{D}\), by hypothesis. We conclude that
\(F_{\mathsf{S}}\left( \mathbf{d} \right) = F\left( \mathbf{d} \right)\)
for all \(\mathbf{d} \in \mathcal{D}\). □~Lemma~C.4

Let \(\left( \mathcal{A}, + \right)\) be an abelian group. For any
subset \(\mathcal{S} \subseteq \mathcal{A}\), let
\(\langle\mathcal{S}\rangle\) denote the divisible subgroup generated by
\(\mathcal{S}\)---that is, \(\langle\mathcal{S}\rangle\) is the smallest
divisible subgroup of \(\mathcal{A}\) which contains \(\mathcal{S}\).

Lemma~C.5

Let \(\left( \mathcal{A}, + \right)\) be an abelian group, let
\(\left( \mathcal{R}, + \right)\) be a torsion-free abelian group, and
let \(\left. \phi:\mathcal{A}\longrightarrow\mathcal{R} \right.\) be a
group homomorphism. Let \(\mathcal{S} \subset \mathcal{A}\) . If
\(\phi(s) = 0\) for all \(s \in \mathcal{S}\), then \(\phi(a) = 0\) for
all \(a \in {\langle\mathcal{S}\rangle}\).

Proof

Let \(\mathcal{K} ≔ \left\{ a \in \mathcal{A};\phi(a) = 0 \right\}\).
Then \(\mathcal{S} \subseteq \mathcal{K}\). We must show that
\({\langle\mathcal{S}\rangle} \subseteq \mathcal{K}\). To do this, it
suffices to show that \(\mathcal{K}\) is a divisible subgroup. Clearly,
\(\mathcal{K}\) is a subgroup. To see that \(\mathcal{K}\) is divisible,
let \(a \in \mathcal{A}\) and \(n \in \mathbb{N}\) and suppose
\(n\, a \in \mathcal{K}\). Then we have
\(0 = \phi(n\, a) = n\,\phi(a)\). But \(\mathcal{R}\) is torsion-free,
so if \(n\,\phi(a) = 0\) then \(\phi(a) = 0\). Thus
\(a \in \mathcal{K}\).□~Lemma~C.5

Lemma~C.6

Let \(\left( \mathcal{A}, + \right)\) be an abelian group. Let
\(\mathcal{S}_{1},\mathcal{S}_{2} \subseteq \mathcal{A}\) be arbitrary
subsets, and suppose that
\(\mathcal{S} ≔ \mathcal{S}_{1} \cup \mathcal{S}_{2}\) is additively
closed. Then either
\({\langle\mathcal{S}\rangle} = {\langle\mathcal{S}_{1}\rangle}\), or
\({\langle\mathcal{S}\rangle} = {\langle\mathcal{S}_{2}\rangle}\).

Proof

Suppose
\({\langle\mathcal{S}_{1}\rangle} \neq {\langle\mathcal{S}\rangle}\); we
will show that
\({\langle\mathcal{S}_{2}\rangle} = {\langle\mathcal{S}\rangle}\).
Clearly,
\({\langle\mathcal{S}_{1}\rangle} \subseteq {\langle\mathcal{S}\rangle}\).
Thus, if
\({\langle\mathcal{S}\rangle} \neq {\langle\mathcal{S}_{1}\rangle}\),
then
\({\langle\mathcal{S}\rangle} \nsubseteq {\langle\mathcal{S}_{1}\rangle}\),
which means that
\(\mathcal{S} \nsubseteq {\langle\mathcal{S}_{1}\rangle}\). Thus,
\(\mathcal{S} \smallsetminus {\langle\mathcal{S}_{1}\rangle} \neq \varnothing\).
Fix some element
\(s_{2} \in \mathcal{S} \smallsetminus {\langle\mathcal{S}_{1}\rangle}\).
Thus,
\(s_{2} \in \mathcal{S} \smallsetminus \mathcal{S}_{1} \subseteq \mathcal{S}_{2}\).

Claim~C.6.1

\(\mathcal{S}_{1} \subseteq {\langle\mathcal{S}_{2}\rangle}\).

Proof

Let \(s_{1} \in \mathcal{S}_{1}\). Let \(t ≔ s_{1} + s_{2}\). Then
\(t \in \mathcal{S}\), because \(s_{2} \in \mathcal{S}\) and
\(s_{1} \in \mathcal{S}_{1} \subset \mathcal{S}\) and \(\mathcal{S}\) is
additively closed.

First note that \(t \notin \mathcal{S}_{1}\). (Proof: (by contradiction)
Note that \(s_{2} = t - s_{1}\), and \(s_{1} \in \mathcal{S}_{1}\). So
if \(t \in \mathcal{S}_{1}\) also, then
\(s_{2} \in {\langle\mathcal{S}_{1}\rangle}\). But
\(s_{2} \notin {\langle\mathcal{S}_{1}\rangle}\) by definition, so this
is a contradiction. To avoid the contradiction, we cannot have
\(t \in \mathcal{S}_{1}\).)

Thus,
\(t \in \mathcal{S} \smallsetminus \mathcal{S}_{1} \subseteq \mathcal{S}_{2}\).
But \(s_{1} = t - s_{2}\). Thus,
\(s_{1} \in {\langle\mathcal{S}_{2}\rangle}\). This argument works for
all \(s_{1} \in \mathcal{S}_{1}\). The claim follows. ♢~Claim~C.6.1

Now,
\[{\langle\mathcal{S}_{2}\rangle}\underset{( \ast )}{\subseteq}{\langle\mathcal{S}\rangle}\underset{( \dagger )}{=}{\langle\mathcal{S}_{1} \cup \mathcal{S}_{2}\rangle}\underset{(\diamondsuit)}{\subseteq}{\langle{\langle\mathcal{S}_{2}\rangle} \cup \mathcal{S}_{2}\rangle} = {\langle\mathcal{S}_{2}\rangle}\text{,}\]
where \(( \ast )\) is because \(\mathcal{S}_{2} \subseteq \mathcal{S}\),
~\(( \dagger )\) is because
\(\mathcal{S} = \mathcal{S}_{1} \cup \mathcal{S}_{2}\), and
\((\diamondsuit)\) is by Claim~C.6.1. We conclude that
\({\langle\mathcal{S}\rangle} = {\langle\mathcal{S}_{2}\rangle}\), as
desired.□~Lemma~C.6

Lemma~C.7

Let \(\left( \mathcal{A}, + \right)\) be an abelian group. Let
\(\mathcal{S}_{1},\ldots,\mathcal{S}_{N} \subseteq \mathcal{A}\) be
cones, and suppose
\(\mathcal{S} ≔ \mathcal{S}_{1} \cup \cdots \cup \mathcal{S}_{N}\) is
additively closed. Then there exists \(n \in \lbrack 1\ldots N\rbrack\)
such that
\({\langle\mathcal{S}\rangle} = {\langle\mathcal{S}_{n}\rangle}\).

ProofBy Induction on \(N\)

If \(N = 1\), then \(\mathcal{S} = \mathcal{S}_{1}\), and the result is
trivially true.

Now suppose \(N \geq 2\). By induction, suppose that, for any collection
\(\mathcal{S}_{1},\ldots,\mathcal{S}_{N - 1} \subseteq \mathcal{A}\) of
\(N - 1\) cones such that
\(\mathcal{S} ≔ \mathcal{S}_{1} \cup \cdots \cup \mathcal{S}_{N - 1}\)
is additively closed, there exists
\(n \in \lbrack 1\ldots N - 1\rbrack\) such that
\({\langle\mathcal{S}\rangle} = {\langle\mathcal{S}_{n}\rangle}\).

Now let \(\mathcal{S}_{1},\ldots,\mathcal{S}_{N} \subseteq \mathcal{A}\)
be a collection of \(N\) cones such that
\(\mathcal{S} ≔ \mathcal{S}_{1} \cup \cdots \cup \mathcal{S}_{N}\) is
additively closed. We must find some \(n \in \lbrack 1\ldots N\rbrack\)
such that
\({\langle\mathcal{S}\rangle} = {\langle\mathcal{S}_{n}\rangle}\). If
\({\langle\mathcal{S}\rangle} = {\langle\mathcal{S}_{N}\rangle}\), then
we are done. So, suppose
\({\langle\mathcal{S}_{N}\rangle} \subsetneq {\langle\mathcal{S}\rangle}\).
Then \(\langle\mathcal{S}_{N}\rangle\) is a divisible proper subgroup of
\(\mathcal{A}\). Corollary~B.1(b) yields a complete, homogeneous
preorder \(( \succeq )\) on \(\mathcal{A}\) such that
\({\langle\mathcal{S}_{N}\rangle} = \left\{ a \in \mathcal{A};a \approx 0 \right\}\).

Define
\(\mathcal{S}^{\succeq} ≔ \left\{ s \in \mathcal{S};s \succeq 0 \right\}\)
and
\(\mathcal{S}^{\preceq} ≔ \left\{ s \in \mathcal{S};s \preceq 0 \right\}\).
Then \(\mathcal{S} = \mathcal{S}^{\succeq} \cup \mathcal{S}^{\preceq}\),
and \(\mathcal{S}\) is additively closed, so Lemma~C.6 says that either
\({\langle\mathcal{S}\rangle} = {\langle\mathcal{S}^{\succeq}\rangle}\)
or
\({\langle\mathcal{S}\rangle} = {\langle\mathcal{S}^{\preceq}\rangle}\).
Without loss of generality, suppose
\({\langle\mathcal{S}\rangle} = {\langle\mathcal{S}^{\succeq}\rangle}\).

Define
\(\mathcal{S}^{\succ} ≔ \left\{ s \in \mathcal{S};s \succ 0 \right\}\)
and
\(\mathcal{S}^{0} ≔ \left\{ s \in \mathcal{S};s \approx 0 \right\} = \mathcal{S} \cap {\langle\mathcal{S}_{N}\rangle}\).
Then
\(\mathcal{S}^{\succeq} = \mathcal{S}^{\succ} \cup \mathcal{S}^{0}\).
Also \(\mathcal{S}^{\succeq}\) is itself additively closed (because it
is an intersection of two additively closed sets) so Lemma~C.6 says that
either
\({\langle\mathcal{S}^{\succeq}\rangle} = {\langle\mathcal{S}^{\succ}\rangle}\)
or
\({\langle\mathcal{S}^{\succeq}\rangle} = {\langle\mathcal{S}^{0}\rangle}\).
However, by hypothesis,
\({\langle\mathcal{S}^{\succeq}\rangle} = {\langle\mathcal{S}\rangle} \supsetneq {\langle\mathcal{S}_{N}\rangle} = {\langle{\langle\mathcal{S}_{N}\rangle}\rangle} \supseteq {\langle\mathcal{S} \cap {\langle\mathcal{S}_{N}\rangle}\rangle} = {\langle\mathcal{S}^{0}\rangle}\).
Thus,
\({\langle\mathcal{S}^{\succeq}\rangle} \neq {\langle\mathcal{S}^{0}\rangle}\).
Thus,
\({\langle\mathcal{S}^{\succeq}\rangle} = {\langle\mathcal{S}^{\succ}\rangle}\).
But we have supposed
\({\langle\mathcal{S}\rangle} = {\langle\mathcal{S}^{\succeq}\rangle}\),
so we conclude that
\({\langle\mathcal{S}\rangle} = {\langle\mathcal{S}^{\succ}\rangle}\).

For all \(n \in \lbrack 1\ldots N\rbrack\), define
\(\mathcal{S}_{n}^{\succ} ≔ \left\{ s \in \mathcal{S}_{n};s \succ 0 \right\}\).
Then
\(\mathcal{S}^{\succ} = \mathcal{S}_{1}^{\succ} \cup \cdots \cup \mathcal{S}_{N - 1}^{\succ} \cup \mathcal{S}_{N}^{\succ}\),
(because
\(\mathcal{S} = \mathcal{S}_{1} \cup \cdots \cup \mathcal{S}_{N - 1} \cup \mathcal{S}_{N}\)).
But \(\mathcal{S}_{N}^{\succ} = \varnothing\) (because
\(s \approx 0\) for all \(s \in \mathcal{S}_{N}\), by the construction
of the homogeneous preorder \(( \succeq )\)). Thus,
\(\mathcal{S}^{\succ} = \mathcal{S}_{1}^{\succ} \cup \cdots \cup \mathcal{S}_{N - 1}^{\succ}\).
Furthermore, \(\mathcal{S}^{\succ}\) is additively closed and
\(\mathcal{S}_{1}^{\succ},\ldots,\mathcal{S}_{N - 1}^{\succ}\) are all
cones (because \(\mathcal{S}\) is additively closed and
\(\mathcal{S}_{1},\ldots,\mathcal{S}_{N - 1}\) are cones, and the set
\(\left\{ a \in \mathcal{A};a \succ 0 \right\}\) is itself a cone).
Thus, the induction hypothesis yields some
\(n \in \lbrack 1\ldots N - 1\rbrack\) such that
\({\langle\mathcal{S}^{\succ}\rangle} = {\langle\mathcal{S}_{n}^{\succ}\rangle}\).
But then we have:
\[{\langle\mathcal{S}\rangle} = {\langle\mathcal{S}^{\succ}\rangle} = {\langle\mathcal{S}_{n}^{\succ}\rangle} \subseteq {\langle\mathcal{S}_{n}\rangle} \subseteq {\langle\mathcal{S}\rangle}\text{.}\]
We conclude that
\({\langle\mathcal{S}_{n}\rangle} = {\langle\mathcal{S}\rangle}\), as
desired. □~Lemma~C.7

Let \(\left. \pi:\mathcal{X}\longrightarrow\mathcal{X} \right.\) be a
permutation, and let \(x \in \mathcal{X}\). The \(\phi\)-orbit of \(x\)
is the set
\(\mathcal{O}_{x}^{\phi} = \left\{ x,\phi(x),\phi^{2}(x),\phi^{3}(x),\ldots \right\}\).
This set must be finite, because \(\mathcal{X}\) is finite. Thus, there
is some \(n \in \mathbb{N}\) such that \(\phi^{n}(x) = x\). Note that
\(\mathcal{O}_{x}^{\phi}\) is also the \(\phi\)-orbit of any other
element in \(\mathcal{O}_{x}^{\phi}\). If \(x,y \in \mathcal{X}\), then
it is easy to see that either
\(\mathcal{O}_{x}^{\phi} = \mathcal{O}_{y}^{\phi}\), or
\(\mathcal{O}_{x}^{\phi}\) is disjoint from \(\mathcal{O}_{y}^{\phi}\).
Thus, we can partition \(\mathcal{X}\) into disjoint \(\phi\)-orbits
\(\mathcal{O}_{x_{1}}^{\phi},\mathcal{O}_{x_{2}}^{\phi},\ldots,\mathcal{O}_{x_{M}}^{\phi}\)
for some \(x_{1},x_{2},\ldots,x_{M} \in \mathcal{X}\); this is called
the \(\phi\)-orbit partition of \(\mathcal{X}\). Lemma~C.8

Let \(\mathcal{R}\) be a linearly ordered abelian group. Let
\(\mathcal{X}\) be a finite set, let \(\mathcal{V}\) be any set, let
\(\mathcal{D} \subseteq \mathbb{N}^{\langle\mathcal{V}\rangle}\) be a
cone, and let
\(\left. F:\mathcal{D}\rightrightarrows\mathcal{X} \right.\) be a
balance rule determined by a perfect \(\mathcal{R}\)-valued balance
system. Then \(F\) is neutral if and only if \(F\) is a scoring rule
with an \(\mathcal{R}\)-valued neutral score system.

Proof

``\(\Longleftarrow\)'' If the score function of \(F\) is neutral, then
\(F\) itself is neutral.

``\(\Longrightarrow\)'' Suppose \(F\) is a perfect
\(\mathcal{R}\)-valued balance system. If \(F\) is
\(\Pi_{\mathcal{X}}\)-neutral, then Lemma~C.3 says that
\(F = F_{\mathsf{B}}\), where \(\mathsf{B}\) is a neutral,
\(\mathcal{R}\)-valued balance system.

First suppose \(\left| \mathcal{X} \right| = 2\). If
\(\mathcal{X} = \left\{ x,y \right\}\), then
\(\mathsf{B} = \left\{ \mathbf{b}^{x,y},\mathbf{b}^{y,x} \right\}\),
where \(\mathbf{b}^{y,x} = - \mathbf{b}^{x,y}\). Thus, if we define
\(\mathbf{s}^{x} ≔ \mathbf{b}^{x,y}\) and
\(\mathbf{s}^{y} ≔ \mathbf{0}\), then \(F = F_{\mathsf{S}}\).

Now suppose \(\left| \mathcal{X} \right| \geq 3\). There are three
cases: either \(\left| \mathcal{X} \right| \equiv 0\) or
\(\left| \mathcal{X} \right| \equiv 1\) or
\(\left| \mathcal{X} \right| \equiv 2,\operatorname{mod}3\). Refer to
these as Cases 0, 1, and 2. Partition \(\mathcal{X}\) as
\(\mathcal{X} = \mathcal{X}_{0} \sqcup \mathcal{X}_{1} \sqcup \mathcal{X}_{2} \sqcup \cdots \sqcup \mathcal{X}_{N}\),
(for some \(N \geq 0\)), where \(\left| \mathcal{X}_{n} \right| = 3\)
for all \(n \in \lbrack 1\ldots N\rbrack\), and where
\(\mathcal{X}_{0} = \varnothing\) in Case 0,
\(\left| \mathcal{X}_{0} \right| = 4\) in Case 1, and
\(\left| \mathcal{X}_{0} \right| = 5\) in Case 2.

Let
\({\widetilde{\Pi}}_{\mathcal{X}} ≔ \left\{ \widetilde{\pi};\pi \in \Pi_{\mathcal{X}} \right\} \subseteq \Pi_{\mathcal{V}}\).
Then \(\mathcal{D}\) is a \({\widetilde{\Pi}}_{\mathcal{X}}\)-invariant
subset of \(\mathbb{N}^{\langle\mathcal{V}\rangle}\), because \(F\) is
neutral. Let \(\Phi \subset \Pi_{\mathcal{X}}\) be the set of all
permutations whose orbit-partition of \(\mathcal{X}\) is
\(\left\{ \mathcal{X}_{0},\mathcal{X}_{1},\ldots,\mathcal{X}_{N} \right\}\).
All elements of \(\Phi\) have the same order
\(m\).\textsuperscript{17}17\(m = 3\) in Case 0. If \(N = 0\), then
\(m = 4\) in Case 1 and \(m = 5\) in Case 2. If \(N \geq 1\), then
\(m = 12\) in Case 1 and \(m = 15\) in Case 2. Let \(M ≔ m - 1\). For
any \(\phi \in \Phi\) and \(\mathbf{d} \in \mathcal{D}\), define
\(\mathbf{d}^{\phi} ≔ \mathbf{d} + \widetilde{\phi}\left( \mathbf{d} \right) + {\widetilde{\phi}}^{2}\left( \mathbf{d} \right) + \cdots + {\widetilde{\phi}}^{M}\left( \mathbf{d} \right)\);
then \(\mathbf{d}^{\phi} \in \mathcal{D}\) also, because \(\mathcal{D}\)
is \(\widetilde{\phi}\)-invariant and additively closed. Thus,
\(F_{\mathsf{B}}\left( \mathbf{d}^{\phi} \right) \neq \varnothing\).
For all \(n \in \lbrack 0\ldots N\rbrack\), let
\(\mathcal{D}_{n}^{\phi} ≔ \left\{ \mathbf{d} \in \mathcal{D};\mathcal{X}_{n} \subseteq F_{\mathsf{B}}\left( \mathbf{d}^{\phi} \right) \right\}\).
Claim~C.8.1

For any \(\phi \in \Phi\) and any \(n \in \lbrack 0\ldots N\rbrack\),
the set \(\mathcal{D}_{n}^{\phi}\) is a cone.

ProofAdditively Closed

Let \(\mathbf{d}_{1},\mathbf{d}_{2} \in \mathcal{D}_{n}^{\phi}\). We
must show that
\(\mathbf{d}_{1} + \mathbf{d}_{2} \in \mathcal{D}_{n}^{\phi}\). By
hypothesis,
\(\mathcal{X}_{n} \subseteq F_{\mathsf{B}}\left( \mathbf{d}_{1}^{\phi} \right)\)
and
\(\mathcal{X}_{n} \subseteq F_{\mathsf{B}}\left( \mathbf{d}_{1}^{\phi} \right)\).
Thus, reinforcement implies that
\(\mathcal{X}_{n} \subseteq F_{\mathsf{B}}\left( \mathbf{d}_{1}^{\phi} + \mathbf{d}_{2}^{\phi} \right)\).
But
\(\mathbf{d}_{1}^{\phi} + \mathbf{d}_{2}^{\phi} = \left( \mathbf{d}_{1} + \mathbf{d}_{2} \right)^{\phi}\);
thus
\(\mathcal{X}_{n} \subseteq F_{\mathsf{B}}\left( \left( \mathbf{d}_{1} + \mathbf{d}_{2} \right)^{\phi} \right)\),
which means
\(\mathbf{d}_{1} + \mathbf{d}_{2} \in \mathcal{D}_{n}^{\phi}\).

(Divisible)~Let
\(\mathbf{z} \in \mathbb{Z}^{\langle\mathcal{V}\rangle}\) and
\(m \in \mathbb{N}\), and suppose
\(m\,\mathbf{z} \in \mathcal{D}_{n}^{\phi}\). We must show that
\(\mathbf{z} \in \mathcal{D}_{n}^{\phi}\). By hypothesis,
\(\mathcal{X}_{n} \subseteq F_{\mathsf{B}}\left( \left( m\,\mathbf{z} \right)^{\phi} \right)\).
But \(\left( m\,\mathbf{z} \right)^{\phi} = m\,\mathbf{z}^{\phi}\).
Thus, \(\mathbf{z}^{\phi} \in \mathcal{D}\) (because \(\mathcal{D}\) is
a cone), and
\(F_{\mathsf{B}}\left( m\,\mathbf{z}^{\phi} \right) = F_{\mathsf{B}}\left( \mathbf{z}^{\phi} \right)\),
by reinforcement. Thus,
\(\mathcal{X}_{n} \subseteq F_{\mathsf{B}}\left( \mathbf{z}^{\phi} \right)\).
Thus, \(\mathbf{z} \in \mathcal{D}_{n}^{\phi}\). ♢Claim~C.8.1

Claim~C.8.2

For any \(\phi \in \Phi\), we have
\(\mathcal{D} = \mathcal{D}_{0}^{\phi} \cup \mathcal{D}_{1}^{\phi} \cup \cdots \cup \mathcal{D}_{N}^{\phi}\).

Proof

For any \(\mathbf{d} \in \mathcal{D}\), we have
\(\widetilde{\phi}\left( \mathbf{d}^{\phi} \right) = \mathbf{d}^{\phi}\).
Thus neutrality implies that
\(\phi\left\lbrack F_{\mathsf{B}}\left( \mathbf{d}^{\phi} \right) \right\rbrack = F_{\mathsf{B}}\left( \mathbf{d}^{\phi} \right)\).
Thus, \(F_{\mathsf{B}}\left( \mathbf{d}^{\phi} \right)\) must be a union
of \(\phi\)-orbits. Thus,
\(F_{\mathsf{B}}\left( \mathbf{d}^{\phi} \right)\) contains
\(\mathcal{X}_{n}\) for some \(n \in \lbrack 0\ldots N\rbrack\). Thus,
\(\mathbf{d} \in \mathcal{D}_{n}^{\phi}\) for some
\(n \in \lbrack 0\ldots N\rbrack\). ♢Claim~C.8.2

Claim~C.8.3

Suppose there exists some \(\phi \in \Phi\) and some
\(n \in \lbrack 1\ldots N\rbrack\) such that
\({\langle\mathcal{D}_{n}^{\phi}\rangle} = {\langle\mathcal{D}\rangle}\)
. Let \(\mathcal{X}_{n} = \left\{ x,y,z \right\}\), where
\(\phi(x) = z,\phi(y) = x\), and \(\phi(z) = y\) . Then
\(\mathbf{b}^{x,y}\left( \mathbf{d} \right) + \mathbf{b}^{y,z}\left( \mathbf{d} \right) = \mathbf{b}^{x,z}\left( \mathbf{d} \right)\)
for all \(\mathbf{d} \in \mathcal{D}\).

Proof

For any \(\mathbf{d} \in \mathcal{D}_{n}^{\phi}\), we have
\(\mathcal{X}_{n} \subseteq F_{\mathsf{B}}\left( \mathbf{d}^{\phi} \right)\).
Then we must have
\[0 = \mathbf{b}^{x,y}\left( \mathbf{d}^{\phi} \right) = \mathbf{b}^{x,y}\left( \mathbf{d} + \widetilde{\phi}\left( \mathbf{d} \right) + {\widetilde{\phi}}^{2}\left( \mathbf{d} \right) + \cdots + {\widetilde{\phi}}^{M}\left( \mathbf{d} \right) \right)\underset{(\diamondsuit)}{=}\left( \mathbf{b}^{x,y} + \mathbf{b}^{x,y}\widetilde{\phi} + \mathbf{b}^{x,y}{\widetilde{\phi}}^{2} + \cdots + \mathbf{b}^{x,y}{\widetilde{\phi}}^{M} \right)\left( \mathbf{d} \right)\underset{( \ast )}{=}\left( \mathbf{b}^{x,y} + \mathbf{b}^{\phi^{- 1}{(x)},\phi^{- 1}{(y)}} + \cdots + \mathbf{b}^{\phi^{- M}{(x)},\phi^{- M}{(y)}} \right)\left( \mathbf{d} \right) = \frac{M + 1}{3}\left( \mathbf{b}^{x,y} + \mathbf{b}^{y,z} + \mathbf{b}^{z,x} \right)\left( \mathbf{d} \right)\text{,}\]
where \((\diamondsuit)\) is by Lemma~C.2(a), and \(( \ast )\) is because
\(\mathsf{B}\) is neutral. But
\({\langle\mathcal{D}_{n}^{\phi}\rangle} = {\langle\mathcal{D}\rangle}\),
so Lemma~C.5 then implies that
\(\mathbf{b}^{x,y}\left( \mathbf{d} \right) + \mathbf{b}^{y,z}\left( \mathbf{d} \right) + \mathbf{b}^{z,x}\left( \mathbf{d} \right) = 0\)
for all \(\mathbf{d} \in \mathcal{D}\). In other words:
\(\mathbf{b}^{x,y}\left( \mathbf{d} \right) + \mathbf{b}^{y,z}\left( \mathbf{d} \right) = - \mathbf{b}^{z,x}\left( \mathbf{d} \right) = \mathbf{b}^{x,z}\left( \mathbf{d} \right)\),
as desired. ♢Claim~C.8.3

Claim~C.8.4

Suppose there exists no \(\phi \in \Phi\) and
\(n \in \lbrack 1\ldots N\rbrack\) such that
\({\langle\mathcal{D}_{n}^{\phi}\rangle} = {\langle\mathcal{D}\rangle}\)
. Suppose that either
\(\mathcal{X}_{0} = \left\{ x,y,z,s \right\}\)~(in~Case~1), or
\(\mathcal{X}_{0} = \left\{ x,y,z,s,t \right\}\)~(in~Case~2). Then
\(\mathbf{b}^{x,y}\left( \mathbf{d} \right) + \mathbf{b}^{y,z}\left( \mathbf{d} \right) = \mathbf{b}^{x,z}\left( \mathbf{d} \right)\)
for all \(\mathbf{d} \in \mathcal{D}\).

Proof

For all \(\phi \in \Phi\), if there is no
\(n \in \lbrack 1\ldots N\rbrack\) such that
\({\langle\mathcal{D}_{n}^{\phi}\rangle} = {\langle\mathcal{D}\rangle}\),
then Claim~C.8.1, Claim~C.8.2 and Lemma~C.7 imply that
\({\langle\mathcal{D}_{0}^{\phi}\rangle} = {\langle\mathcal{D}\rangle}\).

(Case~1)~Let \(\phi \in \Phi\), and suppose
\(\phi(y) = x,\phi(z) = y,\phi(s) = z\), and \(\phi(x) = s\). Then by an
argument similar to Claim~C.8.3, we conclude that
(C.9)\[\left( \mathbf{b}^{x,y} + \mathbf{b}^{y,z} + \mathbf{b}^{z,s} + \mathbf{b}^{s,x} \right)\left( \mathbf{d} \right) = 0\text{,}\]
for all \(\mathbf{d} \in \mathcal{D}\). By choosing other permutations
in \(\Phi\), we can similarly derive the following five equations. For
all \(\mathbf{d} \in \mathcal{D}\),
(C.10)\[\left( \mathbf{b}^{x,y} + \mathbf{b}^{y,s} + \mathbf{b}^{s,z} + \mathbf{b}^{z,x} \right)\left( \mathbf{d} \right) = 0\text{,}\](C.11)\[\left( \mathbf{b}^{y,z} + \mathbf{b}^{z,x} + \mathbf{b}^{x,s} + \mathbf{b}^{s,y} \right)\left( \mathbf{d} \right) = 0\text{,}\](C.12)\[\left( \mathbf{b}^{y,z} + \mathbf{b}^{z,s} + \mathbf{b}^{s,x} + \mathbf{b}^{x,y} \right)\left( \mathbf{d} \right) = 0\text{,}\](C.13)\[\left( \mathbf{b}^{z,x} + \mathbf{b}^{x,y} + \mathbf{b}^{y,s} + \mathbf{b}^{s,z} \right)\left( \mathbf{d} \right) = 0\text{,}\](C.14)\[\text{and}\quad\left( \mathbf{b}^{z,x} + \mathbf{b}^{x,s} + \mathbf{b}^{s,y} + \mathbf{b}^{y,z} \right)\left( \mathbf{d} \right) = 0\text{.}\]
Recall that \(\mathbf{b}^{y,s} = - \mathbf{b}^{s,y}\), and
\(\mathbf{b}^{z,s} = - \mathbf{b}^{s,z}\), etc. Thus, by adding Eqs.
(C.9), (C.10), (C.11), (C.12), (C.13), (C.14) together and canceling, we
obtain
\(4\left( \mathbf{b}^{x,y} + \mathbf{b}^{y,z} + \mathbf{b}^{z,x} \right)\left( \mathbf{d} \right) = 0\),
for all \(\mathbf{d} \in \mathcal{D}\). Thus,
\(\mathbf{b}^{x,y}\left( \mathbf{d} \right) + \mathbf{b}^{y,z}\left( \mathbf{d} \right) = - \mathbf{b}^{z,x}\left( \mathbf{d} \right) = \mathbf{b}^{x,z}\left( \mathbf{d} \right)\),
as desired.

(Case~2) Let \(\phi \in \Phi\), and suppose
\(\phi(y) = x,\phi(z) = y,\phi(s) = z,\phi(t) = s\), and
\(\phi(x) = t\). Then by an argument similar to Claim~C.8.3, we conclude
that
(C.15)\[\left( \mathbf{b}^{x,y} + \mathbf{b}^{y,z} + \mathbf{b}^{z,s} + \mathbf{b}^{s,t} + \mathbf{b}^{t,x} \right)\left( \mathbf{d} \right) = 0\text{,}\]
for all \(\mathbf{d} \in \mathcal{D}\). By choosing other permutations
in \(\Phi\), we can similarly derive the following five equations. For
all \(\mathbf{d} \in \mathcal{D}\),
(C.16)\[\left( \mathbf{b}^{x,y} + \mathbf{b}^{y,z} + \mathbf{b}^{z,t} + \mathbf{b}^{t,s} + \mathbf{b}^{s,x} \right)\left( \mathbf{d} \right) = 0\text{,}\](C.17)\[\left( \mathbf{b}^{y,z} + \mathbf{b}^{z,x} + \mathbf{b}^{x,s} + \mathbf{b}^{s,t} + \mathbf{b}^{t,y} \right)\left( \mathbf{d} \right) = 0\text{,}\](C.18)\[\left( \mathbf{b}^{y,z} + \mathbf{b}^{z,x} + \mathbf{b}^{x,t} + \mathbf{b}^{t,s} + \mathbf{b}^{s,y} \right)\left( \mathbf{d} \right) = 0\text{,}\](C.19)\[\left( \mathbf{b}^{z,x} + \mathbf{b}^{x,y} + \mathbf{b}^{y,s} + \mathbf{b}^{s,t} + \mathbf{b}^{t,z} \right)\left( \mathbf{d} \right) = 0\text{,}\](C.20)\[\text{and}\quad\left( \mathbf{b}^{z,x} + \mathbf{b}^{x,y} + \mathbf{b}^{y,t} + \mathbf{b}^{t,s} + \mathbf{b}^{s,z} \right)\left( \mathbf{d} \right) = 0\text{.}\]
By adding Eqs. (C.15), (C.16), (C.17), (C.18), (C.19), (C.20) together
and canceling, we obtain
\(4\left( \mathbf{b}^{x,y} + \mathbf{b}^{y,z} + \mathbf{b}^{z,x} \right)\left( \mathbf{d} \right) = 0\),
for all \(\mathbf{d} \in \mathcal{D}\), Thus,
\(\mathbf{b}^{x,y}\left( \mathbf{d} \right) + \mathbf{b}^{y,z}\left( \mathbf{d} \right) = \mathbf{b}^{x,z}\left( \mathbf{d} \right)\),
as desired. ♢~Claim~C.8.4

In Case 0, we have \(\mathcal{D}_{0} = \varnothing\), so
\(\mathcal{D} = \mathcal{D}_{1} \cup \cdots \cup \mathcal{D}_{N}\).
Thus, Claim~C.8.1, Claim~C.8.2 and Lemma~C.7 yields some
\(n \in \lbrack 1\ldots N\rbrack\) such that
\({\langle\mathcal{D}_{n}\rangle} = {\langle\mathcal{D}\rangle}\). Thus,
Claim~C.8.3 applies. In Cases 1 or 2, either the hypothesis of
Claim~C.8.3 holds or the hypothesis of Claim~C.8.4 holds. In any of
these cases, we obtain some distinct elements \(x,y,z \in \mathcal{X}\)
such that
(C.21)\[\mathbf{b}^{x,y}\left( \mathbf{d} \right) + \mathbf{b}^{y,z}\left( \mathbf{d} \right) = \mathbf{b}^{x,z}\left( \mathbf{d} \right)\text{,}\quad\text{for~all~}\mathbf{d} \in \mathcal{D}\text{.}\]

Claim~C.8.5

For any \(x',y',z' \in \mathcal{X}\), and all
\(\mathbf{d} \in \mathcal{D}\), we have
\(\mathbf{b}^{x',y'}\left( \mathbf{d} \right) + \mathbf{b}^{y',z'}\left( \mathbf{d} \right) = \mathbf{b}^{x',z'}\left( \mathbf{d} \right)\).

Proof

There exists some \(\alpha \in \Pi_{\mathcal{X}}\) such that
\(\alpha\left( x' \right) = x\), \(\alpha\left( y' \right) = y\), and
\(\alpha\left( z' \right) = z\). Thus,
\[\mathbf{b}^{x',y'}\left( \mathbf{d} \right) + \mathbf{b}^{y',z'}\left( \mathbf{d} \right)\underset{(\diamondsuit)}{=}\left( \mathbf{b}^{x,y}\widetilde{\alpha} \right)\left( \mathbf{d} \right) + \left( \mathbf{b}^{y,z}\widetilde{\alpha} \right)\left( \mathbf{d} \right)\underset{( \dagger )}{=}\left( \mathbf{b}^{x,y} + \mathbf{b}^{y,z} \right)\left( \widetilde{\alpha}\left( \mathbf{d} \right) \right)\underset{( \ast )}{=}\mathbf{b}^{x,z}\left( \widetilde{\alpha}\left( \mathbf{d} \right) \right)\underset{(\diamondsuit)}{=}\mathbf{b}^{x',z'}\left( \mathbf{d} \right)\text{,}\]
as desired. Here, \(( \ast )\) is by Eq. (C.21) (because
\(\widetilde{\alpha}\left( \mathbf{d} \right) \in \mathcal{D}\) because
\(\mathcal{D}\) is \({\widetilde{\Pi}}_{\mathcal{X}}\)-invariant), and
\(( \dagger )\) is by Lemma~C.2(a), while both \((\diamondsuit)\) are
because \(\mathsf{B}\) is neutral. ♢~Claim~C.8.5

Now, Claim~C.8.5 and Lemma~C.4 imply that \(F_{\mathsf{B}}\) is a
scoring rule. Finally, if \(F_{\mathsf{B}}\) is neutral, then
Proposition~1 says that \(F_{\mathsf{B}}\) has a neutral score function.
□~Lemma~C.8

Proof of Theorem~2

``\(\Longleftarrow\)'' It is easy to check that any scoring rule
satisfies reinforcement. If the score system of \(F\) is neutral, then
\(F\) itself is neutral.

``\(\Longrightarrow\)'' If \(F\) satisfies reinforcement, then Theorem~1
says \(F\) is a balance rule. If \(F\) is neutral, then Lemma~C.8 says
\(F\) is scoring rule with a neutral score system.□~Theorem~2

1Smith (1973) calls this condition `separability', while Young, 1974b,
Young, 1975 and~Saari (1990) call it `consistency'. Dhillon (1998) and
Dhillon and Mertens (1999) refer to an analogous property as `extended
Pareto', while Gilboa and Schmeidler (2003) call it `combination'.2Smith
(1973) and Young (1974b) consider rules which produce social preference
relations as output, whereas Young (1975) considers rules which produce
subsets of social alternatives as output; these two frameworks yield two
slightly different definitions of `reinforcement' and `scoring
rule'.3See~Young, 1974a, Young, 1975 and Nitzan and Rubinstein
(1981).4See Young and Levenglick (1978).5See Richelson (1978),
Morkelyunas (1982), Ching (1996), and~Yeh (2008).6Zwicker (2008) has
shown that such scoring rules can also be interpreted as `mean proximity
rules': each element of \(\mathcal{V}\) and \(\mathcal{X}\) is
represented as a vector in \(\mathbb{R}^{N}\), and the social choice is
the element of \(\mathcal{X}\) which is closest to the vector average of
the signals of the voters.7See Fishburn (1978), Morkelyunas (1981), and
Alós-Ferrer (2006).8Smith (1973) also considered linearly ordered
abelian groups.9To be precise: for any
\(\mathbf{r} = \left( r_{1},\ldots,r_{N} \right)\) and
\(\mathbf{r}' ≔ \left( r_{1}',\ldots,r_{N}' \right)\) in
\(\mathbb{R}^{N}\), we have \(\mathbf{r} \gg \mathbf{r}'\) if there is
some \(M \in \lbrack 1\ldots N\rbrack\) such that \(r_{n} = r_{n}'\) for
all \(n < M\), while \(r_{M} > r_{M}'\).10Thus, if
\(\mathcal{R} = \mathbb{R}\), then
\(\mathbf{r}\left( \mathbf{n} \right) = \mathbf{r}\bullet\mathbf{n}\),
where `\(\bullet\)' is the inner product operation on
\(\mathbb{R}^{\mathcal{V}}\).11This example also appears in Myerson
(1995, Section 4, p. 65).12This is slightly weaker than the usual
definition of ``reinforcement'', because we do not require the domain
\(\mathcal{D}\) to be fully closed under addition---we only require it
to be weakly additive. Of course, in almost all of the ``natural''
examples, \(\mathcal{D}\)~will be closed under addition. But we allow
the possibility that \(\mathcal{D}\) is not additively closed, in order
to accommodate certain interesting examples, such as the ``Condorcet''
balance rule of Example~3 below.13This is closely related to one of the
main insights of Saari (1990), although that paper deals mainly with
weak reinforcement. I thank the referee for pointing out this
connection.14A similar remark applies to a balance rule
\(\left. F_{\mathsf{B}}:\mathcal{D}\rightrightarrows\mathcal{X} \right.\).
There may exist other profiles
\(\mathbf{n} \in \mathbb{N}^{\langle\mathcal{V}\rangle} \smallsetminus \mathcal{D}\)
which satisfy the nontriviality condition (1), so that
\(F_{\mathsf{B}}\left( \mathbf{n} \right) \neq \varnothing\).
But the fact that \(F_{\mathsf{B}}\) is well-defined outside
\(\mathcal{D}\) does not imply that society is normatively compelled to
apply \(F_{\mathsf{B}}\) outside \(\mathcal{D}\).15I thank the referee
for suggesting this simplified version of the proof.16This object is
often called the positive cone. However, we are reserving the term cone
for a set which is divisible. Hence the ungainly term
`conoid'.17\(m = 3\) in Case 0. If \(N = 0\), then \(m = 4\) in Case 1
and \(m = 5\) in Case 2. If \(N \geq 1\), then \(m = 12\) in Case 1 and
\(m = 15\) in Case 2.

\end{document}
